% Created 2025-08-21 Thu 22:55
% Intended LaTeX compiler: pdflatex
\documentclass[11pt]{article}
\usepackage[utf8]{inputenc}
\usepackage[T1]{fontenc}
\usepackage{graphicx}
\usepackage{longtable}
\usepackage{wrapfig}
\usepackage{rotating}
\usepackage[normalem]{ulem}
\usepackage{amsmath}
\usepackage{amssymb}
\usepackage{capt-of}
\usepackage{hyperref}
\author{LiquidZulu}
\date{\today}
\title{Anatomy of the Celebritarian}
\hypersetup{
 pdfauthor={LiquidZulu},
 pdftitle={Anatomy of the Celebritarian},
 pdfkeywords={},
 pdfsubject={},
 pdfcreator={Emacs 30.1 (Org mode 9.7.29)}, 
 pdflang={English}}
\begin{document}

\maketitle
\tableofcontents

Intro: Clips of smith et al playing on static-ey TV, switching between channels on cuts, zooms out to massive wall of clips, TVs flash off to form the title ``Anatomy of the Celebritarian'', sombre violin music playing

EN: use edits for transitions between major sections; e.g. supercut of trump being an authoritarian for the trump section.
\begin{itemize}
\item if there is a section on living in ancapistan in one's head, use Danny Duchamp's rick roll
\item maybe ``real women vote for trump song'' for Trump section; or some edit song that goes hard
\begin{itemize}
\item could be good for Angela
\end{itemize}
\end{itemize}
\section{Cargo Cults}
\label{sec:org0e8ebfd}
Among the indigenous tribes of the South Pacific there exists a peculiar phenomenon. In the late 19th century, the islanders began constructing facsimiles of airstrips using twigs and performing military drills with bamboo shoots in the place of guns.

Previously, these tribes had witnessed a great war between nations possessing technology they could not conceive of, and they saw that supplies would be brought to strange places they could not understand the true purpose of. Lacking any context for these great flying machines, they would attempt to entice them to bring gifts by building monuments in favour of whatever gods must have sent them.

These tribes are called cargo cults. Their operating premise is that if they simply build an area which looks like the airports they have seen, that they too will be able to receive the riches bestowed by the sky. That is, they have not understood the fundamental reason behind the activities they saw others engaged in; so they do not know that those activities were not the reason for success.

A similar mentality has gripped the libertarian movement; incapable or unwilling to understand the greats, most have resigned to simply imitating their movements, without understanding the purpose behind those movements. If this poison is not removed, we are doomed to fall into a world of communism. This is a poison that has taken hold of some of the most popular advocates of liberty. This is the anatomy of the celebritarian.

{[}intro edit]
\section{Pragmatism, Old and New}
\label{sec:orgb9bd6cf}
\subsection{The Philosophy of Pragmatism}
\label{sec:org1937166}
In the wake of Trump's tariffs there were hordes of economists pointing out that this would spell economic ruin for the United States. A common response to this was to tell the economists that they were too married to abstract theory, and that there is a crisis in this moment that must be dealt with. These types have it that we shouldn't get bogged down in theoretical arguments or general principles, we must be ``practical'' and address the problem of the day. The attitude we are to adopt is of the manner: try and see in each and every case.

This line of argument comes from the philosophy of pragmatism. Of course, most of the man-on-the-street type advocates of pragmatism, including those in prominent positions within various political movements, have absolutely no clue about the implications of the philosophy they are endorsing. So I think it is worthwhile to go over what exactly pragmatism claims, and what it implies.

The pragmatists come onto the philosophical scene in a time where the possibility for knowledge about the world has been under attack from every direction. They lamented how philosophy had retreated into Noumenal worlds, and Platonic Forms, and Dialectic processes; completely disconnected from man's life. Thus they sought to bring philosophy back down to Earth, so to speak, and make it \emph{work} for the every day man.

Accordingly, they held that the purpose of theory is action, that the purpose of philosophy should be to guide man in every day situations. Now, this is not simply the claim that having the correct theory in place will allow you to make the correct decisions; rather their claim is that the very standard of whether the theory is correct in the first place is whether it works. In a sentence: pragmatism holds that the true is a species of the good.

The pragmatists have developed legion arguments for this theory of truth coming from just about every direction you can imagine in philosophy, but for the purposes of this video it will be sufficient to give just one.

Basically: take Descartes' argument of Cartesian Doubt. That what it means to say that you know something with certainty is to say that you couldn't be wrong about that thing. But man is supposed to be a fallible being, meaning he \emph{can} be wrong about things. So how can a being which can be wrong about things ever be said to be in a state where he can't be wrong about something?

Surely he can't. Thus man must doubt absolutely everything which is dubitable, so he can never know anything with certainty. Of course, Descartes tried to get around this by saying that he knew with certainty that he existed to do the doubting, but subsequent philosophers would throw this out as a completely arbitrary leap.

So, the situation we are in is that on the classical theory of what it means for something to be true, we are left in total epistemic stasis, unable to know a single thing. Thus the pragmatists tell us that we need to abandon this model of what it means to say that something is true. If philosophy is to be a guide, clearly it has failed---this theory of truth is useless. So, how about a theory that says something is true if it works?

The ethical implications of this theory are obvious. If something is true if and to the extent that it works, there can be no unyielding ethical principles. What works today may or may not work tomorrow. We have to wait and see. But, what is the standard of an ethic working or not? Well, the standard of truth in any field, ethics or not, is that it achieves some human purpose. A thing is true because it leads to desirable outcomes to hold it as true. This applies also to logic itself. The laws of logic, on this view, are not self-evident, or axiomatic, or anything of the sort. They are true if and to the extent that they satisfy human goals. Thus the purpose that we are holding as the standard of truth must be pre-logical; that is, it must be just a mere brute desire, or whim. So an ethical claim is true if and insofar as it satisfies one's unjustified and unjustifiable whims. In short, the pragmatist ethics is a gross subjectivism in the order of ``I want it! I want it! I want it!''

One of the major pragmatist thinkers, John Dewey, would say the following regarding his ethics:
\begin{quote}
Men, instead of being proud of accepting and asserting beliefs and ``principles'' on the ground of loyalty, will be as ashamed of that procedure as they would now be to confess their assent to a scientific theory out of reverence for Newton or Helmholz or whomever, without regard to evidence.

If one stops to consider the matter, is there not something strange in the fact that men should consider loyalty to ``laws,'' principles, standards, ideals to be an inherent virtue, accounted unto them for righteousness? It is as if they were making up for some secret sense of weakness by rigidity and intensity of insistent attachment. A moral law, like a law in physics, is not something to swear by and stick to at all hazards [\ldots{}].\footnote{John Dewey, \emph{The Quest for Certainty}, p. 278}
\end{quote}

I think that is about as clear as one can get in outright rejecting the concept of a principled approach to ethical questions.

Of course, there are many objections to pragmatism, not least is that it is a form of the primacy of consciousness: the thesis that consciousness precedes existence. They tell us that the true is that which works; and a thing working means that it achieves certain purposes, purposes held by a consciousness. So the true is that which satisfies the purposes of a consciousness. But the primacy of consciousness is itself a stolen concept fallacy. To be conscious means to be conscious \emph{of something}. The concept ``consciousness'' relies upon the \emph{prior} concept ``existence,'' it simply cannot float on its own.

Due to this most fundamental error committed by pragmatism, it ends up being a wholly parasitic philosophy. In ethics, it is simply not possible to be a consistent pragmatist right from the offset. Pragmatism simply takes as given that you already come into it with a code of values provided by some other school, and then it tells you that given those values here's how you can tell what works to achieve them or not.

Perhaps I desire money, and so the pragmatist sets out to tell me that I should follow whichever activity would ``work'' to yield the highest income. But, why would I desire money in the first place? I only desire money because I have a certain code of values, that tells me money is good because it allows me to obtain things that I should want to obtain. That is, I must have already had some basic, perhaps implicit, ethical principles before I could ever follow the pragmatist's advice. Pragmatism goes absolutely nowhere without these primordial desires that it feasts upon; so the pragmatist ethics is out on this point alone.

The parasitic nature of pragmatism extends down into their epistemology too. Pragmatism tells us that saying a proposition is true \emph{means} that it leads to certain consequences and those consequences are desirable. Take some proposition A. ``A is true'' means that A leads to B and B is desirable. So if I am to know that A is true, I must know that A leads to B and that B is desirable. But how do you establish the proposition that A leads to B? Well, as we know, a proposition is true if it has desirable consequences, so ``A leads to B'' is true if ``A leads to B'' leads to C and C is desirable. We arrive at an infinite regress. The pragmatist is forced to throw up his arms and declare that when he says A leads to B that this just simply is the case. He doesn't need to justify \emph{that} pragmatically as well. In other words, the pragmatist relies on the correspondence theory of truth to establish the higher order claims like A leading to B.
\subsection{Pragmatism and the LP}
\label{sec:org84498ba}
In spite of the utter failure of pragmatism to even do what it was set out to achieve, relying on having a pre-existing guide rather than guiding man on its own, it has remained as a powerful influence on much of the libertarian movement.

For years the LP was under the control of the pragmatist caucus. Their stated goal was seeking ``practical, pragmatic solutions to our nation's problems.''\footnote{They don't seem to actually have a website, only a Facebook page, and I don't have Facebook so I can't access it. I pulled that quote from their bio on \url{https://lpcolorado.org/caucuses/}. Of course, that website is itself a buggy mess that barely stays online, so you can read the bio on the Internet Archive: \url{https://web.archive.org/web/20250603130014/https://lpcolorado.org/caucuses/}} Of course, we know that this means a rejection of hardline principle in favour of what ``works.'' The standard of working that these prags adopted, was whatever generated the most acceptance by regime actors in Washington luncheons. This is what made the prags feel very important and professional; that was their immutable desire and Categorical Imperative. This yielded a distinctly leftist, regime-oriented Libertarian Party. It was under their control that Raytheon lobbyist Bill Weld would be granted the position of Vice Presidential candidate.

Later, in a debate with one of the faces of pragmatism in the LP, former Chair Nick Sarwark, Dave Smith would utterly rebuke the pragmatist strategy:
\begin{quote}
I think that this is a very important topic that gets right at the core of what we are all doing here. Like, what are we doing as libertarians? What is our reason for being? And what is the plan here? [\ldots{}] I became a libertarian about 11 years ago, I was a child of the Ron Paul movement, and like many people here when I was introduced to these ideas they changed my life. [\ldots{}] When these ideas click with you, I was like, ``I'm going to be a soldier in the liberty movement for the rest of my life.'' [\ldots{}] I heard Nick recently, I was watching a video of his when he was doing a debate for the LP chairmanship, and he said, ``as the chair of the LP, the party is paramount.'' I don't look at things that way, [\ldots{}] I don't care about the Libertarian Party. I don't care about it at all. [\ldots{}] What's paramount to me is liberty. I care about the liberty movement, I care about the philosophy of libertarianism, the most beautiful, brilliant political philosophy ever devised by man.\footnote{ReasonTV, \emph{What is the Ideal Stratgy for the Libertarian Party? A Soho Forum Debate}, \url{https://www.youtube.com/watch?v=adou3a9g040}, t. 02:14}

{[}\ldots{}]

The philosophy of libertarianism is the philosophy of peace, it's the philosophy of prosperity, [\ldots{}] it's the philosophy of civilisation. I mean, what better describes civilised behaviour than the NAP and respect for private property rights?\footnote{ReasonTV, \emph{What is the Ideal Stratgy for the Libertarian Party? A Soho Forum Debate}, \url{https://www.youtube.com/watch?v=adou3a9g040}, t. 03:50}

{[}\ldots{}]

The only question that matters is to libertarians as far as national candidates goes is: who can convert more people?\footnote{ReasonTV, \emph{What is the Ideal Stratgy for the Libertarian Party? A Soho Forum Debate}, \url{https://www.youtube.com/watch?v=adou3a9g040}, t. 08:00}

{[}\ldots{}]

I don't care what you think of him: it was the great Ron Paul, he was the guy who converted people. Every one of you knows this. [\ldots{}] What did he do? He was courageous. He was principled.\footnote{ReasonTV, \emph{What is the Ideal Stratgy for the Libertarian Party? A Soho Forum Debate}, \url{https://www.youtube.com/watch?v=adou3a9g040}, t. 08:30}
\end{quote}

So, per Dave's argument here, it is not the case that what ``works best'' for the Libertarian Party is the correct standard by which we should judge a strategy. Rather, we must be first and foremost adherents of the \emph{philosophy} of libertarianism. We must stick to the absolute moral principle that is the Non-Aggression Principle. And further, that the purpose of political strategy is convincing people. The only way to do this is to engage in the Ron Paul model: to be principled and courageous. And so fundamentally, the question burning at the heart of the radical vs pragmatist divide is which approach is capable not just of getting people to call themselves libertarians, but to actually convince them \emph{of libertarianism}.

Pragmatism simply \emph{cannot} by its nature convert people \emph{to libertarianism}. Pragmatism is the opposite of principles, so its proliferation could never lead to the spread of any principled stance, including libertarianism. Dave is entirely correct in this argument, that the pragmatist strategy is doomed to failure, and that it is wrong by its own standards because it is incapable of actually working.

This was the impetus behind the formation of the Mises Caucus. This caucus was to bring about a new radical wave within the Libertarian Party; and drop the pragmatism of old. I will allow Dave to explicate this mission:
\begin{quote}
Everybody here knows that there are a lot of lies told about who we are, so I want to just make it unambiguous, and be very clear about who we are. [\ldots{}] This is who we are: we are the Ron Paul army.\footnote{LP Mises Caucus, \emph{Dave Smith - Take Human Action Bash 2021}, \url{https://www.youtube.com/watch?v=hrHFjy4HhrU}, t. 07:25}

{[}\ldots{}]

We're not secretive about what we're trying to do here with the Mises Caucus. [\ldots{}] We want to bring the liberty movement into the LP; we want to re-create the Ron Paul revolution; and we want to speak for the American people and \emph{demand} our freedom. That is the whole mission.\footnote{LP Mises Caucus, \emph{Dave Smith - Take Human Action Bash 2021}, \url{https://www.youtube.com/watch?v=hrHFjy4HhrU}, t. 09:30}
\end{quote}

This message of radicalism over pragmatism; principle over convenience; confidence over cowardice; would resonate with the libertarian vanguard who had for so long lost all hope that the liberty movement would ever align with the libertarian philosophy again. The Mises Caucus would achieve a landslide victory in the 2022 Libertarian National Convention, in what was called the ``Reno Reset.'' The party was no longer under the control of pragmatists, and was poised to venture out in a new radical direction.
\subsection{The Rise of the Neo-Prag}
\label{sec:org925ec24}
Unfortunately, even in the early days of the Mises Caucus there were warning signs that not everything would go according to plan. There was, as I pointed out at the time,\footnote{LiquidZulu, \emph{These People are Destroying Libertarianism}, \url{https://www.youtube.com/watch?v=064I02yeLhg}} a subversive infection that had gripped the Mises Caucus. An infection of neo-prags.

Where the prags of old were willing to sacrifice principles to appease the establishment, these neo-prags were willing to do the same but for the anti-establishment; both at the expense of truth and ethics. The alternative we were to be presented with was thus between right-pragmatism or left-pragmatism---it was tacitly accepted by many that our strategy must revolve around achieving momentary culture war wins, our only option is whether we throw our lot in with the leftist side of the culture war, or the rightist side. After all, these culture war talking points are the issue of the day; we can't allow principle to get in the way of doing what works to solve it!

This infection has, at present, taken complete control over the Mises Caucus. Any hint of radicalism among the leadership and public activists has been completely stamped out. Every major voice in the caucus has performed a complete 180 from the previous radical messaging. Now we are imbibed by a constant chorus telling us that we must be practical, that we must abandon our principles, that we must stop living in anarcho-capitalism in our heads.

As Dave says with respect to libertarians who support open borders:
\begin{quote}
I think that there is something about being so married to ideology that you will support something that is so obviously right in front of your face going to be a disaster for so many people [\ldots{}]\footnote{misesmedia, \emph{Does Libertarianism Require Support for Open Borders?}, \url{https://www.youtube.com/watch?v=1EZ6-Pn9j88}, t. 04:02}
\end{quote}

So, the standard of the true is whether it has positive or negative outcomes. We are too married to ideology; we need to be more pragmatic and address the problem here and now! Of course, Dave is no longer able to appeal to any principles or ideology by which he could validate that open borders would be a disaster, or any standard for what it even means for something to be disastrous. Again, like all pragmatists, Dave's thesis is utterly parasitic.

He asserts his new pragmatism again in the Biblical Anarchy Podcast:
\begin{quote}
If that was correct by libertarian theory, then I would simply reject libertarianism. If the case was that random homeless drug addicts who are covered in their own faeces and filth should be allowed in public schools, [\ldots{}] then I would just abandon libertarianism.\footnote{The Biblical Anarchy Podcast, \emph{What is the Libertarian Response to Immigration w/Dave Smith}, \url{https://www.youtube.com/watch?v=ghqN\_F2iTdU}, 13:46}
\end{quote}

Again, the counter provided here is not a theoretic one. Dave has attempted to provide theoretic counters to the open borders stance, which I will cover later in this video, but the point is that the argument he is making here is thorougly pragmatist. The fact that certain icky conclusions drop out of the theory is used as proof that the theory is false; that is, because the theory doesn't work by the standard of satisfying his whims, the theory is false. This is pragmatism through and through. It is also quite amusing that he will constantly jump to these drug-based examples whilst holding Ross being freed as the justification and purpose of the new Mises Caucus-pragmatism. He of course, wants to have it both ways, as all pragmatists and Machiavellians do. We must support the legality of drugs in one breath, and use their existence to justify violating rights in the next.

Dave would of course go on to endorse Trump as president following promises of freeing Ross and placing a libertarian in the cabinet; and to his credit, he did apologise after Trump turned out to be just as much of a Tyrant in his second term as he was in his first. But, let's check the reason Dave gave for this support in the first place:
\begin{quote}
I do understand where the strategy that Angela kind of embarked on here in the libertarian party, and a lot of us libertarians, I know for myself, and probably for a lot of other people who were supporting Donald Trump: this was very much a different strategy for all of us. [\ldots{}] I would just kind of implore a lot of libertarians to think about the fact that the world has changed. Truly drastically the world has changed over the last few years. There's a lot of things you could point to of what really changed the world, I particularly think COVID, that's the thing I would really point to [\ldots{}] but the world has changed, and in a situation like that you should at least be open to new strategies. Rather than having this, what seems to me, quasi-religious allergy to even suggesting that maybe just sitting in our corner and preaching our message is not the best use of this party.\footnote{Dave Smith, \emph{Did I sell out my principles by supporting Trump?} (Twitter Space), \url{https://x.com/ComicDaveSmith/status/1860865805006590365}, t. 2:30}
\end{quote}

So, the world has changed and our principles must change with it. What works today might not work tomorrow. We need to be practical and take this situation as it presents itself. I trust that I needn't explicate further the pragmatist basis of this argument and why it is therefore false.
\section{The Correct Strategy}
\label{sec:org85c0613}
\subsection{Isaiah's Job}
\label{sec:org5331b33}
Dave does bring up an interesting point here, that we should be open to any new strategy by the fact that it is new, rather than adhering to the strategy of purely sitting in our corner and preaching the radical message. Now, I'm not sure which period he is dreaming up in which the Libertarian Party was in any way engaged in radical messaging; but regardless, what he is doing here is denouncing the strategy of the man whom he calls ``the greatest living American hero,''\footnote{Dave Smith, \emph{Part of the Problem Exclusive \#22}} Ron Paul. Namely, the strategy explicated by Albert Jay Nock in \emph{Isaiah's Job}.\footnote{Albert Jay Nock, \emph{Isaiah's Job}, \url{https://mises.org/mises-daily/isaiahs-job}}

The article covers a story from the Bible of the prophet Isaiah. God told the prophet that civlisation would fall and there would be mass destruction because the people had fallen into sin. He told Isaiah to go out and preach His message, even though ``The official class and their intelligentsia will turn their noses at you and the masses will not even listen. They will all keep on in their own ways until they carry everything down to destruction.''

Isaiah asked the Lord what the point of such an endeavour would be if the people wouldn't listen. He responded, ``Ah, you do not get the point. There is a Remnant there that you know nothing about. They are obscure, unorganised, inarticulate, each one rubbing along as best he can. They need to be encouraged and braced up because when everything has gone completely to the dogs, they are the ones who will come back and build up a new society; and meanwhile, your preaching will reassure them and keep them hanging on. Your job is to take care of the Remnant, so be off now and set about it.''

Nock draws from this story the most powerful strategy there is. It is the strategy that Ron Paul would base his campaign on. It is the strategy that I have devoted my entire career to. The strategy is that you simply go out, stand on the step, and preach your message, day after day. Don't water it down for the masses, don't shy away from the unpopular truths that must be told. Stand up, grow some balls, and say it with your chest.

Most people will laugh at you. They will call you names. They will say you are a maniac, or an extremist, or an enemy of everything society holds dear. Indeed you are, you must embrace this fact, because in every audience of one-hundred, whilst ninety-nine boo, there will be one who listens---one who is able to listen and understand what you are saying. This rational one in the sea of the unthinking egregore is a member of the remnant. He will join you and together you might preach the message again the next day. Soon you will have a small cadre of supporters. Then you will have an underground movement. Then, in time, even countries will fall under the weight of an idea whose time has come.

This is the strategy that Dave Smith has previously endorsed with all of his might, and now writes off entirely. This is the strategy that has taken us from a handful of people in Murray Rothbard's living room to the country of Argentina. This is the strategy that you must follow if you are to have any hope of success.

The neo-prags have taken up an entirely different strategy: one of being a sort of ``ideological salesman.'' Their argument is that effective libertarian messaging must ``meet people where they are''---namely that we dull down the sharp edges and try to mould the theory to fit whatever culture warrior one is trying to herd into the libertarian tent. This strategy necessitates that one disregard principle.

There are simply two opposing theses here: either Isaiah has the correct strategy, and the theory has primacy over messaging, or the salesman is right and messaging has primacy over the theory. It is possible for people to pussyfoot around and not accept either thesis to its logical conclusions, but those people are simply inconsistent. Any salesman who is not engaged in adulterating the theory to best appeal to the masses is not taking his role seriously.
\subsection{The Philosophical Theory of History}
\label{sec:org604feb3}
The reason why the \emph{Isaiah's Job} strategy works is because it is based in the correct interpretation of what moves a society. Namely, that the thing which has the power to cause the changes in history and determines which societies will inhabit it; is not great men, or dialectic processes, or psychological conditions, but the philosophy of the day.

This is Leonard Peikoff's philosophical theory of history. The idea is that the type of society which obtains at any given point is due to the choices of the men who make it up. Those choices are based upon the premises that those men hold. Because we are looking for the sorts of premises which permeate an entire society; we are looking for the most fundamental set of beliefs that a person can hold. That is precisely the set of beliefs that philosophy tackles.

Namely, if a man believes that knowledge can be gained only by appealing to the mystic divinations of temple priests or Keynesian economists then he will find himself weak and dependent. If the philosophy of the day teaches man that his only moral worth can be found in doing his duty and that this demands he not live for himself, then he will be primed to obey the orders of any tyrant who comes along.

This is the key insight for our topic here. If men broadly accept a philosophy that implies a statist politics, this will practically result in the average man having less and less control over his own life, and more and more dependence upon the government. Through every inevitable crisis that results from such a system, the man will feel ever less in control and ever less able to shape his own life and achieve wellbeing, and he will be accordingly driven deeper into the well of dependence upon the ones who rule him.

All the while the intellectuals spend their time degrading the individual, telling him that his mind is helpless, he cannot know anything, thought is untrustworthy, there are no answers, there are no absolutes, that he must share his wealth, that his race is evil, and that his only goal in life should be to see how fully he can destroy himself for the collective.

Again, this is why the \emph{Isaiah's Job} model of strategy is correct. The primary goal of any ideological strategy must be to change the philosophy to one that adheres to and implies said ideology. To bring about a libertarian world, we have to push the libertarian \emph{philosophy}; we need a this-wordly, individualist philosophy to be accepted as the philosophy of the day.

You don't do that by watering down your message and trying to reach normies on their own terms. The defining characteristic of the normie is that he follows the norms of the day. We need to \emph{provide} those norms, not take them as given and try to warp them to suit libertarian ends. That strategy necessarily fails because it consists in herding over and creating fiat libertarians who never agreed with our beliefs in the first place. That is, it doesn't actually change the philosophy of the day; it just gives it a yellow paintjob.
\subsection{Libertarian Leninism}
\label{sec:org1f3ab63}
Of course, the insights provided by \emph{Isaiah's Job} are only a starting point for a full theory of strategy. By far the best elaboration upon this starting point is found in Springtime of Nations' series on ``Libertarian Leninism,''\footnote{See: Springtime of Nations, \emph{Libertarian Leninism Part 1 - Opportunism and Adventurism}, \url{https://www.youtube.com/watch?v=a5L3DNF8uVw}; and: Springtime of Nations, \emph{Libertarian Leninism Part 2 - What is to be done?}, \url{https://www.youtube.com/watch?v=1tToYQQgpyk}} which you should absolutely watch, as I will borrow heavily from their analysis here. You should take note that when they say libertarian Leninism, they mean that they are taking the strategic insights which Lenin pioneered, and applying them to the philosophy of libertarianism. The people behind Springtime of Nations are not the so-called ``libertarian'' socialists we hear so much about. They are hardcore anarcho-capitalists.

It is, after all, unquestionable that many ideologies we would consider utterly repugnant have had far greater success than libertarianism. This is due, of course, to the philosophers setting them up for slam-dunks, but that does not tell us why those \emph{specific} ideologies rose to prominence. That is, the fact that Hume paved the way for Kant who paved the way for Hegel does not tell us that Marxism-Leninism would be the specific Hegelian variety to rise up in Soviet Russia. All that the philosophical theory of history tells us is that some manner of Hegelian collectivism was due to take over---it does not tell us which \emph{specific type} of Hegelian collectivism would win out. Certainly, Lenin was faced with numerous opponents, and he controlled a small minority in Russia; and yet he took the entire country in relatively little time. So how did he do it?

Murray Rothbard, the greatest hero of libertarianism, saw the answer in what he called ideological entrepreneurship. Just as in economic entrepreneurship, your strategy must be responsive to the current situation of the day; and as Mises would point out, timing is everything. Now, this is not to be construed as a pragmatist argument. I am not here saying that based on the current going-ons you must mould your principles to deal with them---I am saying that the specific actions required to advance those absolute principles depend upon your circumstances.

In pursuit of a radicalist strategy, Lenin identified two opposite but equally dangerous extremes. On the one hand were the adventurists, who demanded rash and immediate direct action, along with a complete refusal of established political channels. And on the other hand were the opportunists who had a tendency towards half-way measures and compromise, without a strong commitment to long-term goals.

The adventurists are the all-or-nothing sorts. A libertarian adventurist would completely reject a reduction in the income tax on the grounds that it doesn't completely abolish the state with a single swipe of a pen. They are less relevant for this video, but they are also in error; you will see these types frequenting agorist circles, or bemoaning the election of Milei. That is, the error they commit is that they conflate defending against aggression as being the same thing as aggression. If Milei gets into office and with every move cuts away at the state; thus protecting the property of innocents, they will say that he is evil and worthless because he didn't take it all down. Whatever legitimate criticisms one might have of Milei, this is not one of them. Stopping a suicide bomber from killing 10 out of the 100 people he tried to kill does not mean that you are responsible for the deaths of those 90 that did die.

The opportunists, on the other hand, are the pragmatists I have been speaking about. Whichever side of the culture war they find themselves on, they absolutely adore compromising on any principle so long as it gains them brownie points with their chosen sect. The end-goal of a world which adheres to the non-aggression principle is completely dropped for ``the most important election of our lifetime'' or whatever.

Lenin was somewhat unique in that he was not an adherent of either extreme, and his strategic flexibility would allow his movement to survive the failed 1905 revolution in Russia. That is, radical tactics were proper \emph{during} the revolution, but the later years of revolutionary collapse and reaction were times for caution and retreat. Lenin faced the adventurist ``all-or-nothing'' error in the form of Alexander Bogdanov, who called for a futile armed uprising at a time where they simply did not have the resources to succeed. Such an uprising would surely spell the end of any hopes Lenin had of becoming the ruler of a communist dictatorship.

Lenin would speak to this directly:
\begin{quote}
During the Revolution we learned to 'speak french', i.e to [\ldots{}] raise the energy of the direct struggle of the masses and extend its scope. Now, in this time of stagnation, reaction and disintegration, we must learn to 'speak German', i.e. to work slowly until things revive, systematically, steadily [\ldots{}] winning inch by inch.\footnote{Lenin, \emph{The Liquidation of Liquidationism}, \url{https://www.marxists.org/archive/lenin/works/1909/jul/11.htm}}
\end{quote}

After resisting Bogdanov's adventurism, Lenin also avoided the pitfall of opportunism in 1917 when the Bolsheviks were part of the big-tent February Government which had deposed the Tsar. The bolsheviks were a minority in this centrist liberal-dominated government, led by Alexander Kerensky. The turning point would come from General Kornilov's attempt to overthrow the government and restore the monarchy.

This pivotal moment presented both danger and opportunity for the Bolsheviks. On the one hand, the adventurists were tempted to continue their previous tactics of all-out opposition to Kerensky's regime, and not participate in the battle. On the other hand, the opportunists wanted to fight unconditionally. This unprincipled action might well have demoralised the bolshevik millitants and undercut the larger strategic goal of a bolshevik revolution.

Lenin's solution would avoid both fatal errors. He chose to come to the aid of Kerensky against Kornilov's forces, but not for free. Lenin demanded radical and painful concessions from Kerensky, including: arming the workers, bringing bolshevik troops to the fore, and legalising peasant takeovers of landed estates. From his weakened position there was little Kerensky could do about this, and it led to the bolsheviks being accepted as the de facto leaders of the forces opposed to Kornilov, despite not being a majority.

All that Lenin required was to have enough power to flip the war in Kerensky's favour, it is not necessary that he be able to win it outright. If the bolshevik troops are sufficient to turn the tide in favour of Kerensky then he is faced with a win-or-lose scenario. The only possibility for victory is for the bolsheviks to support him.

The lesson learned from this scenario can be easily applied to the case of national elections. Even though the Libertarian Party makes up a fraction of the votes, if it accrues enough power to have a potential spoiler effect, then the LP effectively dictates which party will win. For instance, perhaps without the LP running in some area, the Republicans would get 51\% of the vote, with the Democrats getting 49\%; but if the LP does run then the Republicans fall behind the dems. In such a situation, the libertarian candidate is able to dictate whatever terms he desires. He can demand extreme concessions from both parties, and whichever one accepts the terms can be made to win.

This is the electoral strategy that the LP should be focused on: building up substantial enough and disciplined enough voter blocs to use in spoiling elections. Meanwhile the presidential nominee's job should be to use the stage afforded to him to do Isaiah's job. We are nowhere near a position that the LP can actually win the presidency, so the goal on the national level should have absolutely nothing to do with gaining votes---you already have the disciplined blocs that can affect real change on the local level. The benefit of the presidential nominee above all others is that he is provided a massive speaker through which to speak the message every damn day.

Certainly, this should not be taken as the Categorical Imperative of strategy; if there is a good reason to use the spoiler effect on the national level, then go for it, but this will tend to be a lot more rare: the guy who has been governor of Astawan for the past decade is going to be a lot more long-term oriented and incentivised to actually follow through on concessions made to the LP than the guy who can only ever be president twice. As has been seen with the most recent Trump election, it really doesn't matter to him at all whether he follows through on putting a libertarian in his cabinet, or cutting the deficit, or anything of the sort. He can never run again, so there is nothing holding him to his promises.

Republicans are at the moment more closely aligned with libertarianism for the most part. There is simply no equivalent of a Ron Paul, Thomas Massie, or Justin Amash among the Democrats. If liberatarians are the bolsheviks in this situation, the Republicans would be like Kerensky's liberals, leaving the Democrats as Kornilov's forces. This means we have a certain strategic alignment with the Republicans, but we should absolutely not uncritically support them. They must \emph{earn} our support by implementing libertarian policies.
\section{The Wynand Strategy, and Its Inevitable Failure}
\label{sec:org1f4814c}
\subsection{Gail Wynand and Pragmatism}
\label{sec:org7ab8e11}
Gail Wynand is a central character in Ayn Rand's \emph{The Fountainhead}. He is the owner of several tabloid newspapers which dish out the most surface-level, bottom of the barrel stories imaginable. This yields him a massive audience, which he believes he is in total control of:
\begin{quote}
Power, Dominique. The only thing I ever wanted. To know that there's not a man living whom I can't force to do---anything. Anything I choose. The man I couldn't break would destroy me. But I've spent years finding out how safe I am. They say I have no sense of honor, I've missed something in life. Well, I haven't missed very much, have I? The thing I've missed---it doesn't exist.\footnote{Ayn Rand, \emph{The Fountainhead}}
\end{quote}

Because Wynand is able to smear men like our protagonist, Howard Roark, with lies and half-truths, he believes he can make his readers believe anything. When the time comes for Wynand to help Roark, he finds that he had no real power in the first place. No matter what he does his writers keep leaving and his papers go unsold. In pandering to this bottom-of-the-barrel audience he has become their slave, not their master. It was the audience that was controlling him the whole time. He is forced to publish only views that appeal to the average, not the great.

This same failure is sure to befall the neo-prags who have taken to do whatever it takes to court a right-wing audience. In their pragmatist scrambling for momentary culture war wins, they have not shifted rightists in a libertarian direction---they have only cultivated a mass of followers who hate liberty. This new audience is not under the control of the pragmatist pundit any more than Wynand's audience was under his control.

Just as Wynand generates only fiat followers, the pragmatist can, at best, generate only fiat libertarians. That is: libertarians in name only. These soundbites that you will hear from this sort of pragmatist serve only to reframe one's already existing philosophy in libertarian terms. This is not enough to make someone a true libertarian. What is required for that is actually changing the philosophy they hold, which means supplanting it with a new philosophy.

It is the defining mark of the culture-warrior that he is incapable of modelling anybodies motivations, or analysing any event, outside of the lens of the culture war. So if, for instance, you decide not to support Donald Trump, the culture-warriors in the Mises Caucus decide that you must be a lolbertarian prag, or some variety of gay race communist. It has already been ordained from on high that the based and red-pilled libertarians are supposed to support Trump, and that's what it takes to ``live in reality.'' There is simply no possible motivation they can conceive of other than that. The idea of a principled individual who votes not based upon the swelling of culture war sentiments, but on ethics and reason, simply does not enter their mind. The ``most principled motherfucker you know'' label has already been homesteaded by Dave, after all; so to be principled means to be a rightist.

It is because of this culture war lens these people have grafted onto their eyeballs that makes them lump all criticism under the same banner. The Mises Caucus have been under attack from the very beginning by what might be termed the ``loser brigade'': an association of leftist pragmatists and federal infiltrators who were previously in control of the Libertarian Party for many years. Wherever any criticism is levied against the neo-prags, they are sure to characterise it all as being under this same lolbertarian, leftist, fed banner.

This tendency has made them blind to what they are doing. They have been alienating the radicals whilst insulting those radicals and refusing any manner of debate. This is why, in his Twitter space regarding his endorsement of Trump, Dave allows Jeremiah Harding to prattle on for minutes uninterrupted, or Rebekah Massie to explain her entire, irrelevant backstory; but the second that MonarchoHoppean gets on the mic he is constantly interrupted. If it is made clear that there is a radical libertarian wing still out there, it becomes a lot harder to dismiss criticism.

I think it is worth explicitly tabulating the interruptions here, so you can fully grasp what I am getting at. Harding spoke first, uninterrupted for over six full minutes, in which I counted him making eight separate points. He would later go on to refute Jacob Winograd's suggestion that he was in any way the best pick for a strong libertarian dissenter to open the space with, stating that he is not 100\% libertarian.

So much for the professional lolcow; Rebekah Massie, stating that she ``is somewhere in the middle'' speaks for roughly five minutes, without even making a single point for the first two minutes of her speech. And I say speech, because much like an audience in some auditorium, none of the neo-prag opponents decide to stop her from talking.

Contrast this with an actual principled, open-borders libertarian, and friend of the channel, MonarchoHoppean. He immediately went into talking about libertarian principles, quoting directly from Hans Hoppe to show what he saw as a departure from Hoppe's political strategy; and he is barely able to speak for more than 10 seconds without someone butting in to argue.

When all is said and done, Dave is then able to claim that he listened to the counterarguments, and was unmoved. The controlled opposition presented by the loser brigade gives him all the ammunition he needs to dismiss any radical libertarian argument out of hand.
\subsection{Dave Smith is Bad at Debating}
\label{sec:orgea35ba1}
Gail Wynand's failure on this point gets at what is perhaps the most fundamental problem with Dave Smith as an advocate for libertarianism. That is, he must restrict his debates to being against regime pundits and leftist infiltrators because he just does not know the philosophy even 1\% as much as he thinks he does. He feels emboldened to prance around and claim that he knows his stuff and that people like me don't, but the truth is revealed the instant that Dave steps out of the hug-box debates he has made a name in dominating.

Dave has said that he refuses to debate me, and that he ``genuinely [doesn't] know who the opponent would be''\footnote{Dave Smith, \emph{Part of the Problem Members Only 22}.} on the topic of his support for immigration restrictions; and that he is ``thorougly unimpressed with the arguments that everyone was making'' in his twitter space on the topic; and that he ``[doesn't] have interest in beating up on weak libertarians.''

Ok, let's take a look and see. Let's analyse this master of debate Dave Smith when he faces an actual debate bro, Andrew Wilson. Let's see if his performance stacks up with the picture he is painting of himself. I will note that I must limit myself in this discussion to only a few choice clips; if you want to ensure that I am tracking the flow correctly, I do highly recommend you bask in the glory that is the entire stream. It should highlight to you just how far astray the celebritarian soundbite model of debate will get you when faced with any competent opponent.

Andrew would start the debate with some very apt foreshadowing of what was to come:
\begin{quote}
I want to let you know the reason I'm having this debate with you is because I believe I can bring you to the dark side of the force. And you're already half-way there. My hope is to convince you to get to the rest of the way there.\footnote{The Crucible, \emph{Andrew Wilson vs Dave Smith: Is Libertarianism better than Christian Populism?}, \url{https://www.youtube.com/watch?v=YqIaiQ-aK\_s}, t. 01:40}
\end{quote}

Indeed, Wilson is entirely correct in his assessment here, and his further assessment of Dave's supposed-radicalism elsewhere:
\begin{quote}
Yeah, I think that the radical libertarian claim: I'm not sure that Dave can really hold onto that for much longer. He's not that radical of a libertarian, he's been moving more staunchly to Christian ethics, and he seems to think that some good central authority is necessary and good.\footnote{The Crucible, \emph{Cenk is a scumbag liar}}
\end{quote}

It will indeed take very little for Dave's soundbitatarianism to utterly crumble under even basic philosophical interrogation by the much more knowledgeable Andrew Wilson.

\begin{quote}
Dave, can you demonstrate for us how you own yourself?

Well, what exactly do you mean by that?

Do you own yourself?

Yes.

Can you demonstrate that?

We are both demonstrating it right now on this show.

Well you're not demonstrating that you own yourself.

What do you mean by ``own?''

Well I think that ownership of a thing is just a social construction; that ownership itself is just a paramater of agreed-upon criteria. [\ldots{}] I own that guitar because we agree that I own it---that's it. I don't think that there is anything which you can point to that would demonstrate that this exists outside of a social construction which we agree to.

Ok. I'm missing your point here.

My point is that I think its a social construction when you say that you own yourself.

I didn't say it's not a social construction.\footnote{The Crucible, \emph{Andrew Wilson vs Dave Smith: Is Libertarianism better than Christian Populism?}, \url{https://www.youtube.com/watch?v=YqIaiQ-aK\_s}, t. 09:54}
\end{quote}

So right here Dave has made a critical mistake. First, Dave has at the root of his entire legal theory self-ownership, and from that he gets everything else. I think this is a poor way to ground libertarian legal theory, which you can see me expound at length on in my course, \emph{The Fundamentals of Libertarian Ethics}; but whether self-ownership is a good grounding is not his specific problem here. Wilson asks him to demonstrate that he owns himself, and instead of Dave coming into this with a definition of ownership and asserting it, he asks Wilson what \emph{he} means by ownership.

This is not a good start and it only gets worse. When Dave re-asks Wilson what he means to own something, he tacitly accepts the definition that ownership is a mere social construction. This is a line of argument that Dave, in his decades of libertarian proselytising, should have encountered and formulated a response to. It is abundantly clear that Dave simply lacks the ability to foresee where Andrew is very clearly taking this conversation. If ownership is a social construction, that social construction can change and the entire ethical system built upon it falls.

Dave's shoddy understanding of the concept of ``ownership'' is made clear a short while later when he tries to run argumentation ethics:
\begin{quote}
I own the guitar because we agree I do you only own yourself because we agree you do?

I don't agree that that's what ownership means, but you said yes and I said yes. So we're both exactly even---

Then I want to know how you own yourself.

Well if you actually look up ownership in a dictionary, [\ldots{}] do I exercise control over myself? Do I do so morally? Yes, I think so. And I actually think that in order for us to even have this conversation, we both have agreed to that principle.

I don't think we have.

It's me vs you: do you own the arguments that you're making?

No, how do you own an argument?

No, when I say that I mean: is it yours?

They came out of my mouth, but how would that establish that I ``own'' them?\footnote{The Crucible, \emph{Andrew Wilson vs Dave Smith: Is Libertarianism better than Christian Populism?}, \url{https://www.youtube.com/watch?v=YqIaiQ-aK\_s}, t. 11:39}
\end{quote}

So, again, Dave is clearly way out of his depth here. First, he tries to run the line that they're both ``even'' in the debate because they both said yes to the social-construction definition of ownership. But this is not true in the slightest. The stated point of contention is: ``is libertarianism better than Christian populism?'' For Andrew to win this debate, he only has to show that it is not the case that libertarianism is better. That is, if both stances are equal, Andrew wins by default. The burden is on Dave to show that libertarianism is \emph{better} than Christian populism; not that they are equal in terms of justifiability.

Second, Dave engages in one of the classic blunders of debate. When he is pressed, once again, for his demonstration that he owns himself, a perfect opportunity to bring forth the Rothbardian conflict-avoidance theory of property that he should know like the back of his hand, Dave instead appeals to dictionary definitions. Dave is in the realm of a highly complex, contextual subfield of philosophy. The Merriam-Webster theory of property is not going to save him here. Merriam-Webster simply cannot provide all of the necessary context or justifications for a specific definition within the context of a heterodox legal theory.

Third, Dave then tries his bastard form of argumentation ethics by asking whether Wilson ``owns'' his arguments. This is utter nonsense, that is rightly called out by Wilson. Even by Dave's Merriam-Webster theory of property, it doesn't make sense to say that you own an argument, and it doubly so gets him nowhere closer to a real argumentation ethics justification of the non-aggression principle or self-ownership.

\begin{quote}
Look, define ownership.

I think that the definition of ``ownership'' would be a social construct in which you and I, via agreement criteria, decide that X thing becomes property that is assigned to me.

By that definition would you own yourself?

I would say by the definition of that social construction, I could say that you own yourself.

Ok.

But that doesn't mean anything, let me explain why: because if the social construct changes, and we decide that you can own a person, Dave.

Ok, what if we change the construction and say I own your guitar and you?

Right! That's exactly right! What if we change the social construction of ownership to mean you can own a person; does that mean that people no longer own themselves?\footnote{The Crucible, \emph{Andrew Wilson vs Dave Smith: Is Libertarianism better than Christian Populism?}, \url{https://www.youtube.com/watch?v=YqIaiQ-aK\_s}, t. 12:38}
\end{quote}

The very fact that Dave thought he was going to get one over on Andrew by tuning this challenge back at him speaks to Dave's utter incompetence. He did not understand, whilst it was happening, that this challenge demolishes his entire libertarian framework. Turning it back at Andrew and saying ``well what if it means I own you now?'' establishes \emph{Andrews} point that the social construct can change, and cedes the point that ownership is a social construct. This is not even a wild hypothetical, there have been periods in history where people thought that men could own other men. To establish his point, Dave would need a \emph{natural} theory of law; \emph{not} this positivism that he has swallowed.

Dave continues:
\begin{quote}
The difference here is that I would say that ownership would imply that you are justified in that claim.

How are you justified in that claim?

How am I justified in the claim of owning myself?

Yeah, you just agreed that it was a social construction, so if we agree that it is a social construction then we can own other people via social construction, that would be justified.\footnote{The Crucible, \emph{Andrew Wilson vs Dave Smith: Is Libertarianism better than Christian Populism?}, \url{https://www.youtube.com/watch?v=YqIaiQ-aK\_s}, t. 13:34}
\end{quote}

So, here Dave is again faced with the problem he has had right from the offset of the debate. He lacks any standard of justification. He has not once mentioned the non-aggression principle and used that as his standard; he just keeps rolling over and accepting whatever Wilson gives him. Wilson now has on the flow that the justification of ownership, because ownership has been established as a social construct, is that it is socially accepted. That \emph{is} the justification of owning something as far as the flow is concerned, and Dave fails to counter this with an actual theory of what counts as justification:
\begin{quote}
Well, I don't think so. The major difference there is that as we're having this conversation, and what's the goal here? [\ldots{}] At least to persuade the audience who is listening to us. Because we recognise that they own---let's use the word control because ``own'' is clearly triggering---

I'm not triggered by the word ``own.''

Ok, when I say ``triggering'' I mean that it's gonna start a whole nother debate about semantics, so---

Yeah, but semantics are important in a debate.

So we are both presuming that the people who are listening to this are in control of their own thoughts, and we could potentially persuade them towards our cause. So if we are doing all this, we are recognising in some sense that it is more valuable to persuade people than it is to just enforce our will on them. So in that sense, yes there is some kind of norm that there's a justifiable path where you can persuade people rather than force them.

Sure and I would agree that you can persuade all sorts of people and that would be part of a societal norm and people could agree that that's part of a societal norm---(Dave: right.)---that does not demonstrate that you own yourself. Thus the key principle for the NAP, which is self-ownership, that's where the violation comes in. Do people own themselves: is that a justified position or isn't it?\footnote{The Crucible, \emph{Andrew Wilson vs Dave Smith: Is Libertarianism better than Christian Populism?}, \url{https://www.youtube.com/watch?v=YqIaiQ-aK\_s}, t. 14:03}
\end{quote}

So, Dave has failed again to challenge the social construction point. He gets close to argumentation ethics but fails to cinch it. He has merely established that there is this sense of justifying a proposition, meaning to convince at least the audience of its veracity. Andrew is sure to accept this point openly, and re-frame it as an acceptance of the social norm theory of property. Andrew says ``I would agree that you can persuade all sorts of people and that would be part of a societal norm and people could agree that that's part of a societal norm'' and instead of asserting that he is talking about some manner of objective norm rather than social construct, Dave responds: ``right.'' This is surely as bad of a performance as it is possible for a libertaran to give on this topic. It should be disqualifying to ever present yourself as a libertarian debate expert if you get this far into a discussion on rights whilst ceding completely the argument that ownership is a social construct and never once employing the non-aggression principle as an objective normative standard for rights.

Dave responds to the point of whether people owning themselves is a justified position:
\begin{quote}
I'm saying that yes, I think it is---

Then can you demonstrate it?

Well like what would you mean by that?

Prove that you own yourself.

That I can control myself?\footnote{The Crucible, \emph{Andrew Wilson vs Dave Smith: Is Libertarianism better than Christian Populism?}, \url{https://www.youtube.com/watch?v=YqIaiQ-aK\_s}, t. 15:50}
\end{quote}

Jesus fucking Christ Dave this is not difficult. You should be able to justify the Rothbardian theory of property with two hands tied behind your back, whilst balancing on a unicycle, and having a stick rammed up your ass. How about take Rothbard's argumentum e contrario from \emph{The Ethics of Liberty}; that there are three possible ethics with respect to ownership:
\begin{enumerate}
\item everyone owns themselves;
\item one group owns another, or;
\item everyone owns everyone else.
\end{enumerate}

Rothbard shows the contradiction involved in the other alternatives, leaving us with only the principle of self-ownership. You should not struggle with this, and you should not be the one who is conflating control and ownership. This is day one shit that you are floundering on. You should never have to ask what Andrew means by a justification of self-ownership; this should be automatic after a decade of presenting yourself as \emph{the} libertarian to beat, as \emph{the most consistent motherfucker you know}. I'll tell you what, Dave, I speak to several \emph{way more} consistent people than you on a daily basis. I doubt you would make my top 100. If you want to have this branding as the libertarian debate lord, you should be able to deal with these questions, they just are not that rare or that difficult.

Again, this is not like Dave was minding his business and some Socratic figure emerged from the bushes and started interrogating a position he has never thought deeply about. This is Dave agreeing to a debate with a person who one can find ample evidence of their knowledge of philosophy; and this person then starts interrogating Dave on what is both the topic of debate and something he has spent over a decade studying, and Dave completely flounders. Dave can truly only ever win debates when his opponents are retarded, the moment he comes against a person who is competent in philosophy he is destroyed.

Dave continues floundering for some more time, until Wilson brings the reductios ad absurdum of incest and suicide:
\begin{quote}
Should the state come in and stop your brother from unaliving himself or not, Dave?

I would want them to, yes.

Yes! Would you want the state to stop a man from having sex with his sister? If it would lead to a non-reproductive act? [\ldots{}]

You're talking about adults?

Yeah. [\ldots{}]

Yeah, I guess I would want them to.\footnote{The Crucible, \emph{Andrew Wilson vs Dave Smith: Is Libertarianism better than Christian Populism?}, \url{https://www.youtube.com/watch?v=YqIaiQ-aK\_s}, t. 26:34}

{[}\ldots{}]

Can you justify why two consenting adults who are brother and sister having sex is morally wrong?

Yes.

Then do it.

Ok, because I think it's--for the same reason you would think--I think it's against nature, it's evil, it just seems horrible. I don't know, do you have a better answer for why it's wrong?

Well yes, Divine Command.

Divine Command, Ok so I also say Divine Command.\footnote{The Crucible, \emph{Andrew Wilson vs Dave Smith: Is Libertarianism better than Christian Populism?}, \url{https://www.youtube.com/watch?v=YqIaiQ-aK\_s}, t. 31:43}
\end{quote}

The problem here is that Wilson will obviously have an argument as to why Divine Command establishes the legitimacy of state action to quell degenerate or sinful behaviour. Nobody who is well-versed in debate theory is surprised by where Wilson will take this. Now Dave is in the hole of justifying self-ownership from the starting point of Divine Command. If he doesn't even know what is perhaps Rothbard's most famous argument for self-ownership, I have my doubts that he will be able to do this. Surely Christian populism will fall out of the theory of Divine Command far more neatly than libertarianism.

It is because of Daves unprincipled pragmatism; his refusal to stick by the NAP in situations like incest or suicide where it just produces so many yucky outcomes that he cannot navigate around; that is why he lost this debate. Dave chronically fails to bite any bullets that even a moderately competent debate bro could send his way. He's able to argue against Andrew Cuomo, or Laura Loomer because they are fundamentally unimpressive people, by Dave's own admission. He says outright on the Your Welcome podcast that he has a policy of refusing to debate anybody who knows more about the topic than him:
\begin{quote}
An organisation reached out to me about debating Benny Morris. And I was like: ``no! That's ridiculous! That's just like absurd!'' I've read a lot of books about the Israel-Palestine conflict but he's \emph{written} a lot of books---like I've read \emph{his} books.\footnote{Michael Malice, \emph{``YOUR WELCOME'' with Michael Malice \#315: Dave Smith}, \url{https://www.youtube.com/watch?v=iMR4YbGD2EI}}
\end{quote}

What is extra funny is that Dave attempts to spin this into a dunk on Destiny, saying that Destiny is a person who is more interested in winning debates than being correct in his ideas. But he says this with respect to Destiny \emph{actually} debating Benny Morris. Surely it is the man who refuses to debate people who are more well-read than him who is concerned only with winning debates, rather than the people who go into such debates.

It is worth noting that this podcast was recorded shortly after his debate with Wilson. Perhaps it is a new policy in the wake of his utter failure to win debates with competent people.

\begin{quote}
If we were to have the non-aggression principle, with a shrunken government, are we still going to use the government to aggress against consenting adults who want to have incest?

Well as long as a government exists, it will be used to aggress against non-consenting adults.

Great! Then, yes.

Andrew, you're trying to obsess over this incest argument, Ok; laws against incest. I'll grant you that. Sure. [\ldots{}] It's pretty funny, because we live under the most powerful centralised government in the history of the world. [\ldots{}] Instead the argument we are getting is: but if there was no government, who would stop an adult brother and sister from fucking? That's kind of goofy man.

Saying it's kind of goofy isn't an argument, you just said if we scale the government down, as long as there was still a government you would want them to enforce actions via violating the non-aggression principle.\footnote{The Crucible, \emph{Andrew Wilson vs Dave Smith: Is Libertarianism better than Christian Populism?}, \url{https://www.youtube.com/watch?v=YqIaiQ-aK\_s}, t. 50:07}
\end{quote}

Wilson is entirely correct in his analysis of Dave here. Dave has conceded the entire point of the debate, the non-aggression principle, over and over again. And saying that Andrew's argument is ``goofy'' establishes absolutely nothing.

Dave will do this very frequently. He is perfectly willing to advance the most hardcore reductios ad absurdum when it comes to open borders, but the moment anybody gives him the same treatment he resorts to guffawing over how ``goofy'' the argument is. I call this the ``dipshit Dave gambit.''

That is: when faced with criticism over his misrepresentations of the theory, or over lacking any principles, Dave will often find himself responding with something like, ``I don't know man, it just seems kinda goofy to me,'' and this is supposed to serve in the place of an argument. This is perfectly in line with the Wynand strategy: he is not trying to appeal to the intelligent, principled men. He is trying to herd a gaggle of morons under his control. Something ``sounding goofy'' is perfectly sufficient as an argument for them.

Dave should just be honest. He wants to be in this big arena, where he can grow his following but without getting any significant push-back. He wants to preach to the choir. He wants to always be in a situation where he has a friendly audience and a not-that-great opponent such that he may be glazed and showered with support for how based and red-pilled he is. The \emph{moment} he steps outside of this hug-box, his entire philosophy crumbles.

The only thing Dave is capable of doing is rattling off his pre-prepared soundbites in front of midwits who have no way of responding to them. This involves a confusion between having various applied positions within ethics and having an actual ethical theory. Dave thinks that because he can say that murder is bad, or war is bad, or rape is bad, that that constitutes his ethics. No. That is not what ethics means. Those are \emph{applications} of some ethical theory that he has tried to cherry pick without having the basis from which to derive them. This soundbite-theory of debate is perfectly sufficient when you are dealing with midwits who don't know how to interrogate those soundbites, but it is not even close to enough when you are on the level of \emph{real} debate.

Putting aside cases like incest or suicide, Dave has been so destroyed in this debate that he can't even bite the bullet that fucking conscription is always bad:
\begin{quote}
Dave. The argument that Andrew is making is that duties from the prism of theist morality entail the violation of self ownership, which the NAP banks on. It's a contradiction of 2 worldviews.

Noerr: Basically if I can interpret it, what he's saying is that if you adopt Christian ethics, then there is morality which comes from that that is directly oppositional to the idea of self-ownership.

Ok, so like, give me an example.

Conscription.

Ok, but in what scenario.

In the most charitible situation you can possibly think of.

Ok so if you can say conscription in the most charitible scenario, fine, but then maybe you just have to own the principle of conscription in general? So if we're just gonna argue that it's like: yeah, I think conscription in general is horrifically immoral.

But there's at least certain cases where it's not, right?

I mean maybe, yes.

But that's what he's getting at.

Fine! But if we're moving out of just the realm of pure theory, into the realm of the real world, then conscription has been pretty evil.\footnote{The Crucible, \emph{Andrew Wilson vs Dave Smith: Is Libertarianism better than Christian Populism?}, \url{https://www.youtube.com/watch?v=YqIaiQ-aK\_s}, t. 1:45:04}
\end{quote}

So, not only does Dave concede the point that conscription, that is enslaving people to fight in a war, can sometimes be good; but he also paints this as the ``theoretical'' standpoint. That in ``theory'' conscription is fine, but ``in the real world'' it has been bad. This is a complete \emph{inversion} of the truth. The libertarian theory says that conscription is always bad, because people own themselves. You know, the very premise of this entire debate. It is on Dave's ``what works?'' pragmatism that you can get insane, evil conclusions like slavery being fine sometimes---we just have to wait and see if it works, after all!

This is about the weakest rebuttal possible for any libertarian to bring forth when asked whether conscription is immoral. This is a question that 99\% of baby libertarians could get right, and Dave gets it wrong. There is simply no sense in which he is a libertarian, or even close to being the most consistent libertarian out there, as he loves to claim. I could print out the usernames of every member of my discord and throw a dart at them; the person whose name was pierced by that dart would be able to provide a superior argument regarding the libertarian position on conscription than this.

In fact, let's test this. I went on my stream, selected a random member from the live chat, and asked them whether they were willing to bite the bullet on conscription.

Everyone in the chat has proven that they are literally a more consistent motherfucker than Dave Smith. This is really not a difficult bullet to bite or a challenging argument to make. People own themselves, so the state has no right to enslave them.

Finally, at the very end of the stream, a competent proponent of the anarcho-capitalist legal theory, and follower of mine, Kasimir comes on to interrogate Dave a little more:
\begin{quote}
I'm only cutting this off for this reason: I believe the next caller will have a lot to say on this very specific question, because I know him, and I know he's a big proponent of the non-aggression principle. Kasimir, go ahead sir. [\ldots{}]

I'll get right to the meat and potatoes. [\ldots{}] Dave, do you agree that if I want to build a home for myself, I could use a tape measure?

Sure.

Ok. If I take this tape measure and I measure a 2x4 and I know that a 2x4 is like 8 foot long, and then this bloody thing tells me: ``oh no, it's 12 foot.'' [\ldots{}] this thing tells me: ``it \emph{ought} be 12 foot.'' Would that be a good tape measure?

I don't even understand what you are saying. [\ldots{}] The tape measure told you it \emph{ought to be}?

Wilson: If you measured the wood and it said it was X length and the tape measure said it was a different length it would be bad, right?

Sure.

Kas: Yeah, it would be like a useless one, right?

Ok.

So, do we need the tape measure to tell me what ought be the length? No, right? It should tell me what \emph{is} the length.

Sure.

The point I'm trying to make is that: when we go into a legal system, we also don't need it to tell us what we ought do with it, it just tells us what \emph{is} the crime. In fact, if my legal system were to tell me: ``you ought do this'' or ``you ought do that,'' it would be like the tape-measure telling me what the length should be. I don't need it to tell me what the length should be, I need it to tell me the reality. Is it a crime in reality descriptively or not? Is that not true?

No, I mean, listen, I think that's a terrible point, and I think even Andrew would have to admit by his own standards of morality [Andrew recoils at this], [\ldots{}] if the legal system were to tell you that it's a crime to get married and have kids, or that it was a crime to not torture your own kid, you could go: I don't care what it tells me \emph{is} the situation, that ought not be what I believe in---

There's an analogy-breaker here, this tape-measure is not just arbitrary, like these examples, this is based on reality. The 2x4 is actually 8 foot long, [\ldots{}] this is not an arbitrary invention, this is this is an invention that is based on observing reality and measuring it. So if there was a legal system that was based on objective reality, meaning it was an objective legal system, and what it gave you is a true or false result, whether something is a crime, would you want it to do anything else? [\ldots{}]

I think actually, Andrew's initial point, that he was arguing with me takes out what you're saying here: yeah, but you can't get from an is to an ought.

I'm not trying to get to an ought! I'm saying that ``ought'' is bad when I'm talking about a system of law. What a system of law is to tell me is: what is a crime? It shouldn't tell me: oh, don't perform the crime for X reasons. I don't care about that initially, something else can tell me why I shouldn't perform a crime, for example: Christian ethics can tell me why I shouldn't be a criminal; or it could be a whim: I don't want to be a criminal; or another reason: argumentation ethics---

Ok, so the law is you can't believe in Jesus.

Well that law would not be objective. It would not be like the tape measure---

Well that's convenient. [\ldots{}] The law is objectively that you're not allowed to believe in Jesus---

So this is an equivocation---

I can measure it, and that's what the law is.

That's an equivocation, so, when I talk about ``law''---

Maybe I'm missing the point, so just ask the question.

Dave, what is a conflict?

What is a conflict? [he said, utterly confused.]

Yes.

Ok, to views that don't align with each-other.

No, in libertarian terms it means a contradictory action. [\ldots{}] a conflict is contradictory actions, and from there you can build a system where you could define what could be a system of resolving or avoiding conflicts. And then you come up with three possible options: you could have either libertarian law, or objective law, which is: you should avoid conflicts; you could have an answer which says the law of the jungle: you should not avoid conflicts, or you can do whatever you want; and then there could be some mixed law which is what we're currently living under, [\ldots{}]. Now, all of those, reduce to either might makes right or conflict avoidance. So that is an objective standard. If I can examine what you're doing and saying you're initiating conflict, I can call you a criminal. That's objectively true. I don't care whether it should be or shouldn't be---

Ok, fair enough, I got you.

Let's say you say ``why is that true,'' well, because everything else resolves into just conflict. I don't care what you call a conflict initiation, you can call it gentle caress, [\ldots{}] it still falls under a different category than non-aggression. From not initiating conflicts. So that is an objective tape-measure of law. [\ldots{}]

Yeah, I tend to agree with you on that. There's something about the morality of law, and the question of what is right and wrong are in different categories.

It's [legal] for me to evict somebody who is starving to death, but it might be not moral according to some system.

Yes, I agree with you on that.\footnote{The Crucible, \emph{Andrew Wilson vs Dave Smith: Is Libertarianism better than Christian Populism?}, \url{https://www.youtube.com/watch?v=YqIaiQ-aK\_s}, t. 2:43:27}
\end{quote}

So, Kasimir is providing \emph{my argument} of the NAP from a natural law standpoint, and Dave responds with ``I agree with you on that.'' This is, of course, after a debate throughout which Dave was completely incapable of providing his own foundation. So I am to now accept that Dave is using \emph{my} foundation for law? Interesting. So will he debate me about law? Nah, I just don't know my stuff.

Kasimir would shout me out directly as a superior proponent of the legal theory that Dave is trying and failing to approach:
\begin{quote}
I want to shout out great contemporary jurists like LiquidZulu, Stephan Kinsella, and Lukasz Dominiak. And I hope that they listen, perhaps hear this, join together in a podcast so you can actually observe contemporary libertarian jurists discuss questions of law. Possibly separate from questions of morality.\footnote{The Crucible, \emph{Andrew Wilson vs Dave Smith: Is Libertarianism better than Christian Populism?}, \url{https://www.youtube.com/watch?v=YqIaiQ-aK\_s}, t. 2:52:02}
\end{quote}

This would not be the only time that I specifically was called out as a superior proponent of this stuff than Dave, even within this specific stream. Dave's opponent, Andrew Wilson, also called me out earlier in the stream as being better at arguing the argumentation ethics point:
\begin{quote}
By the way, Andrew, do you have a better answer for than me, of what's right and wrong? [\ldots{}]

Usually libertarians will use argumentation ethics, which I know you're a big Hoppe fan---

That's a different thing. I love argumentation ethics---

It is actually a decent way to create justifications around moral oughts [\ldots{}]. I mean Zulu does this all the time, he seems to be pretty good at this, a guy named LiquidZulu, at creating justifications out of argumentation ethics, so I've enjoyed watching those debates [\ldots{}]\footnote{The Crucible, \emph{Andrew Wilson vs Dave Smith: Is Libertarianism better than Christian Populism?}, \url{https://www.youtube.com/watch?v=YqIaiQ-aK\_s}, t. 1:40:00}
\end{quote}

So, how about we return to Dave's claim that he ``doesn't know who the opponent would be'' and that he won't debate me because he is ``afraid of beating up on weak libertarians.'' Look, Dave, if you ever grow the fucking balls to watch this video and you are thus seeing this---I'm a grown up, I can take care of myself, and I can certainly perform vastly better under this basic interrogation that Andrew Wilson was giving to you. It is an absolutely absurd claim to make that I am unworthy of debating you because you are so much better at it than me. Grow up, grow some balls, and accept my challenge.

In this very stream where you completely embarrassed yourself by not being able to answer basic questions any libertarian has faced, I was specifically called out \emph{twice} as a better proponent of this stuff than you. There just is no denying my prowess when you have two separate people telling you, to your face, that I am better at this shit than you are.

You can keep running whatever gambits you want to get out of debating me, but they keep falling one after the other. You can say that you won't debate me because I'm just not good enough on libertarian theory so you would easily crush me, but that has been shown as completely false by this debate.

You can say that you won't debate me because I'm just too irrelevant for you, but that has been shown to be false by my double-citing in the debate you are having as a better proponent than you, and by the fact that my top video easily sits in your top five videos of all time, and I did it without interviewing fucking presidential candidates.

Sure, LiquidZulu will have no impact on the culture, but here's LiquidZulu receiving a shoutout from Shadiversity;\footnote{\url{https://x.com/shadmbrooks/status/1779874039982678186}} but here's LiquidZulu engaged in a multi-video response chain with TIK History;\footnote{\url{https://www.youtube.com/playlist?list=PLNSNgGzaledhruhd0EakCPLkAZ4jRDikh}} but here's LiquidZulu turning young Objectivists into anarchists; but here's LiquidZulu getting more views than the podcast I got this clip from; but here's LiquidZulu debating the founder of the biggest internet movement since Net Neutrality, and showing up in random drama videos on the topic; and here's LiquidZulu living in \emph{your} head rent free.

Dave, you can keep giving out these excuses to not debate me, or you can give whatever the real reason is. Maybe you just don't like me; well, you debated Andrew fucking Cuomo, so I guess you guys are best buds.

Maybe I just fit into that category of people who are far more well-read on this topic than you are and who are therefore banned from debating you. It would just be nice if you would embrace this fact, explain to everyone that you are a coward who prefers preaching to the choir, and then we can all move on.

Because that's what I find particularly insulting about Dave Smith. It's not just that he is bad at expounding these ideas beyond a few carefully selected soundbites---its that he does this whilst painting himself as nothing less than ``the most consistent motherfucker you know.'' There is absolutely no way on Earth that Dave fucking Smith, of all people, is deserving of this title. This is why I have began homesteading it from him. \emph{I} am the most consistent motherfucker you know. Give it a shot, Dave. Ask me any hardcore NAP question and watch me bite the bullet like you are incapable of doing. Watch me tell people to their face that you should not steal a penny to stop the martians from blowing up the entire planet. A limp-wristed, pragmatist coward like Dave Smith doesn't deserve to be in the same fucking room as a title like ``the most consistent motherfucker you know.'' That's something you need to \emph{earn}.

So Dave, if you feel like I have been unfair to you in any way, man the fuck up and debate me. Prove to me that I am inconsistent where you are consistent. Take that title back from me and I will forever sing your praises. I certainly will not be holding my breath.
\subsection{These People are an Embarrassment}
\label{sec:orgdf917a9}
The problem with Dave's performance here is not just that he lost a debate on a technical level or anything. The problem is that he is representing himself as not only a spokesman of libertarianism, but \emph{the best} spokesman of libertarianism.

Dave understands this fully:
\begin{quote}
When I do a debate, I take debate seriously because, from my perspective,  my job is to represent a worldview, and to represent that worldview to the best of my ability.\footnote{Dave Smith, \emph{The Crucible w\&\#47; Andrew Wilson | Part Of The Problem 1134}, \url{https://www.youtube.com/watch?v=NcDAn9WDfV8}, t. 01:23}

{[}\ldots{}]

I agree with that, and I feel the same way often in debates, that like yeah, Ok, I can kind of express this a little bit better than a lot of people who want to express this thing, so let me make sure I do a good job on that.\footnote{Dave Smith, \emph{The Crucible w\&\#47; Andrew Wilson | Part Of The Problem 1134}, \url{https://www.youtube.com/watch?v=NcDAn9WDfV8}, t. 36:35}
\end{quote}

The fact that he has \emph{chosen} to put himself in this position leaves him open to extensive criticism on how he is performing in that role. If you cultivate a career built around being a prominent libertarian who knows all of the facts and can beat anybody in a debate, it reflects on the wider movement when you fail so miserably. I am not interested in criticising everyone who has ever lost a debate before. But in this particular case Dave has branded himself as an expert and has used that expert branding as a justification to refuse to debate me or anybody on my side of these issues. That is a problem.

We must look no further than Destiny's reaction to this debate to see the effect it is having. Again, I will be able to play only a few choice clips here to paint a general picture of the response to this debate from non-libertarians; but I absolutely urge any interested parties in watching Destiny's entire reaction to this stream to see just how obvious Dave's failing is to any non-sympathetic observer. Of course, we had pretty much daily coverage from every libertarian podcast on the face of the fucking planet, going over every detail of Dave's debate with Alex Nowrasteh, but we get relative crickets when it comes to the Andrew Wilson debate. That should tell you something.

\begin{quote}
I heard this was a bloodbath. Dave Smith tried to debate Andrew Wilson and just---if you're going to debate Andrew Wilson you need to be like, up on your philosohpy. You need to have some really strong fundamental understanding of what your core beliefs are. [\ldots{}] and if you're a libertarian you probably don't have any core beliefs because it is one of the most memeiest ideologies on the internet.\footnote{Destiny, \emph{Dave Smith Gets Mad And Resorts To Insults After Getting Stumped In Libertarian Debate}, \url{https://www.youtube.com/watch?v=EEjUp-QdfgA}, t. 00:32}
\end{quote}

So, right off the bat the stakes are clear. Everyone who is even slightly familiar with Andrew Wilson already knows that he is knowledgeable about philosophy. If Dave did a moment of preparation, this would have been made obvious to him and so the audience will expect him to be prepared for that. Destiny has openly asserted that he believes libertarians have no core beliefs, so when Dave fails to disabuse him of this notion the framing is not just that Dave lacks any sort of explicit philosophy, its that \emph{libertarians} do.

\begin{quote}
The reason why this question is so difficult for Dave to answer is that he has absolutely no moral foundation whatsoever. [\ldots{}]\footnote{Destiny, \emph{Dave Smith Gets Mad And Resorts To Insults After Getting Stumped In Libertarian Debate}, \url{https://www.youtube.com/watch?v=EEjUp-QdfgA}, t. 53:27}
\end{quote}

Destiny is correct on this front. One would think that a prominent libertaran would immediately go to the NAP as their framework through which to analyse the legality of acts of incest, but Dave fails this entirely. It is such an obvious place to take the conversation that even Destiny who does not study philosophy and is not even slightly libertaran can see that that is the obvious thing to do:
\begin{quote}
``The key principle for the non-aggression principle which is self-ownership [\ldots{}] do people own themselves: is that a justified position or isn't it?''

If you're a libertarian you absolutely go for the NAP. [\ldots{}]\footnote{Destiny, \emph{Dave Smith Gets Mad And Resorts To Insults After Getting Stumped In Libertarian Debate}, \url{https://www.youtube.com/watch?v=EEjUp-QdfgA}, t. 29:49}

{[}\ldots{}]

``I'm not even understanding what you're saying---''

I know! We all understand, you haven't been understanding for like the last 15 minutes.\footnote{Destiny, \emph{Dave Smith Gets Mad And Resorts To Insults After Getting Stumped In Libertarian Debate}, \url{https://www.youtube.com/watch?v=EEjUp-QdfgA}, t. 33:39}

{[}\ldots{}]

The idea that a libertarian is having the NAP explained to him is probably the single most humiliating thing that is imaginable.\footnote{Destiny, \emph{Dave Smith Gets Mad And Resorts To Insults After Getting Stumped In Libertarian Debate}, \url{https://www.youtube.com/watch?v=EEjUp-QdfgA}, t. 35:24}

{[}\ldots{}]

If you've read anything about libertarianism you would have started here, this is ground zero. And if you think I'm being harsh against Dave, he is the one who says that he is a debater and he markets himself as an intellectual.\footnote{Destiny, \emph{Dave Smith Gets Mad And Resorts To Insults After Getting Stumped In Libertarian Debate}, \url{https://www.youtube.com/watch?v=EEjUp-QdfgA}, t. 35:48}
\end{quote}

If the person who has to google what the non-aggression principle is is able to tell you that that's what you should start talking about, you can in no way paint yourself as a competent proponent of libertarianism. Again, I'm not holding Dave to some unreasonable standard here, I am holding him to the standard set \emph{by his own branding}. Any day one libertarian debate bro does better than Dave here.
\subsection{Dave Smith Failed Libertarianism 101}
\label{sec:org78add42}
\subsubsection{The ``Retroactive Slavery'' Argument}
\label{sec:org59cfe47}
I want to be clear here about what the reason for Dave's loss is. Certainly many knowledgeable people have lost debates before because they get nervous or they just don't know how debate works or how to control the flow of the conversation. This is certainly true of Dave, but it is not the fundamental reason for his defeat. That fundamental reason is that Dave simply lacks any solid theoretic foundation.

Despite having had access to some of the greatest minds in the liberty movement for many years now, it seems that Dave never really bothered progressing beyond the point of bumper-sticker slogans or vague talking points intended to be chanted by a crowd at a protest.

Let's take a look at an attempt Dave has made on a friendly podcast to explain libertarian theory as it pertains to so-called government property:
\begin{quote}
To say that government property is unowned; let's just take this back to libertarianism 101 for a second, because I know there's a lot of people who will say you know that like as you mentioned earlier the response to my position on this will be like: ``hey but libertarianism 101.'' But actually I'm very good at libertarianism 101.

So let's just for a second take it back to that. So there's basically three major foundational beliefs of libertarianism 101 and they are: self- ownership, the non-aggression principle, and private property rights.

Now, most people if you argue with them philosophically--who are not Libertarians--they will actually grant you self-ownership and the non-aggression principle.

{[}\ldots{}]

So there's this connection that libertarians make which I will admit is the trickiest of all of it is that if you think about it just on the surface why is it that libertarians connect self-ownership and the non-aggression principle to property rights? Like why is it that you would go okay because you own yourself how exactly do you transform that into owning land?

Because that does at least on the surface seems like a little bit of like okay there's a leap there right like it's not it's not as if you know even in the Lockean sense or in the more defined Rothbardian sense that you believe you can mix your labor with the Earth and there therefore it becomes an extension of your body you know what I mean like that does seem a little strange?

So what libertarians are really saying is this: we're not saying if you stick your foot on a plot of land that land now becomes part of your body that's not what we're saying but we're saying that okay if there is an unowned piece of land and if we agree that you own yourself and it's wrong to aggress against other people and let's say there's a plot of land that is unowned maybe for example it is the undiscovered frontier or something like that. And you go out there and you build a house on a piece of land. Well if you build that house and you build that so you and your family can live in that house and you put your labor into that, you work and you build this house and you do that you have hurt or robbed from no one, you've just worked and built something and if someone else were to come along and take that house from you then they have in effect retroactively enslaved you.

You worked to build this house and that wasn't a voluntary transaction you weren't working to build it for this guy but he just took that from you so you were in effect his slave for all of those many hours that you labored to build this house, right?

So like that's kind of the essence of where Libertarians go from this idea of self-ownership and non-aggression to property rights [\ldots{}].

Like there's only one answer here where you could say nobody's rights are violated: nobody was retroactively enslaved or enslaved in order to have this, whereas if someone else comes and takes it from you then you've been victimized, right?\footnote{The Biblical Anarchy Podcast, \emph{What is the Libertarian Response to Immigration w/Dave Smith}, \url{https://www.youtube.com/watch?v=ghqN\_F2iTdU}, t. 38:00}
\end{quote}

I regard this ``retroactive slavery'' argument to be about the most primitive, Lockean bullshit imaginable. It is right up there with Dave's use of ``mixing one's labour with the land.'' That has always been an overly-romantic and utterly vague term that has no place in a rational, objective legal theory. Can I measure my labour? Can I pour a cup of labour and mix it up with a cup of flour to make a labour cake? No. Labour simply \emph{is not the root of property}. Labour is something you do \emph{with} your property.

It is only fitting that Dave would justify this primitive, romantic phrase with an equally primitive justification. That what it means to mix your labour with the land is that if some person comes along later and takes that product from you that you are somehow retroactively enslaved? What utter nonsense. It is simply the case that in the moment when that man is building that house he is not possessed by any other person. He is in literal terms not anybodies slave. This fact does not change because somebody comes along and possesses \emph{the house} that he built. Slavery refers to possession of a human, not possession of houses.

At no point has the man who has stolen the house been in possession of the house-builder, or the architect who designed it, or the loggers who provided the wood. All of these men contributed their labour towards the building of the house. Why does Dave have it that only the builder of the house has been enslaved? Is the builder also the one who has been enslaved if instead that builder sold the house to some other person, who then had it stolen from them? Or has that person now been enslaved for the time they spent at their job to afford the house? If they were a slave in that work, was it an illegitimate contract? Should their employer get that money back? Have they still been enslaved if they got the house for free? If the thief kills them before stealing the house, has he enslaved a person who is now dead? Does this mean that an archaeologist taking artefacts from some long-dead pharoah has enslaved that pharoah? Or has he enslaved the slaves of the pharoah?

If retroactive slavery is possible, does this mean I can also retroactively revoke consent? After all, if I build a house with the intention of only me having it and living in it; but then I meet someone and decide to gift them my house---they surely aren't retroactively enslaving me, right? So then I am able to retroactively consent on behalf of my previous self who built the house. That is, my intention at the time doesn't matter.

But perhaps later on I decide that I regret gifting the house and want it back. So can I now retroactively make the retroactive slavery into slavery once more? That is, reach back through time again to revert my intention of the house building back to being only for myself? After all, the intentions in the moment do not matter, so similarly my intentions in the moment of gifting must not matter, so I must be able to change them too. I decide by retroactive fiat that I never intended to gift the house and always intended to keep it for myself. We have always been at war with Eastasia, after all. How Dave hopes to establish trade on a theory of atemporal consent like this is beyond me. And, as it turns out, beyond him.

You can see at this point, I'm sure, just how easy it would be for any moderately competent philosopher or debate-bro to poke one million holes in this theory. This is exactly why Dave avoids these sorts of debates. The gaggle of morons who are paraded in front of him on Piers Morgan to claim that Israel should kill every Muslim in the world are certainly not of the calibre required to find the problems with Dave's philosophy. And that's fine! If Dave wants to just be a soundbitatarian and go on podcasts to preach the soundbites, more power to him. But he should be honest about this. He should certainly \emph{not} be prancing around claiming to be the \emph{most consistent} libertaran there is and failing to present the theory when faced with day one criticisms.

He can't have it both ways. Either Dave is this master debater who knows the theory back to front, in which case he should stop being a pussy and debate me. Or he is a weakling who must hide behind stronger advocates like myself when faced with opponents who have studied philosophy. Whatever the choice, he must be open and honest about his decision, and stick to that decision. We can't have pragmatists like Dave pretending to represent the vanguard any more than we can have Nick Sarwark pretending to be Ron Paul.
\subsubsection{The Primitive Justification as Non-Fundamental}
\label{sec:org9e9fd6f}
Certainly, Dave is hinting at a somewhat valid justification of libertarian theory. Obviously we have to remove all of this retroactive slavery nonsense, but if you properly establish that men own themselves and can come to own other things through homesteading or trade, then you establish the entire corpus of libertarian theory. But this is still erroneous, even when Rothbard tries to do it.

The reason for this is that self-ownership, homesteading, and contract theory are simply not fundamental. This sort of justification is taking a short-cut way up into the legal theory and making those distant conclusions into axioms. The mathematical equivalent would be starting with the Pythagorean theorem and trigonometry and then trying to justify the rest of triangle geometry from there rather than starting with Euclid's axioms and building up to these distant facts about triangles. Certainly, you could justify a whole bunch of lemmas and proofs by starting with Pythagoras and trig; but this would be entirely arbitrary. You would place yourself in a very weak position, philosophically speaking, by standing on multiple floating clouds, rather than by digging down and building a sturdy foundation first.

The proper root of law, which we must use as our foundation, is not that men own themselves, and that they can own other things, and that they can trade things. We must get rid of these ``ands'' if we are to find the fundamental. Ownership itself, whether it be in one's body or in other goods, depends in the first case on the fact of scarcity, giving rise to the possibility of conflicts. This is the primordial fact that law is based on: conflicts (meaning contradictory actions) are possible, so how should we deal with them? It is from this that you recognise three possible types of answer: the non-aggression principle, the law of the jungle, or some manner of mixed law. I have gone over each of these and proven that the non-aggression principle is the only rational solution elsewhere, so I will not belabour this point.

The relevance for the case of Dave and the other weak libertarians who don't actually know any philosophy is that \emph{this is the root}. We need to dig down to the level of understanding what question we are even supposed to ask before we can sail on up into talking about owning yourself or having contracts or anything like that. When one doesn't dig down to the fundamentals, we get situations like Dave advocating government restrictions on incest, or militarising the border, or any of his other numerous errors. This same sloppiness, as an aside, is why the Objectivists so often endorse obvious absurdities like ``intellectual property''---that is, they failed to dig down to the root question and see that the issue is that conflicts are possible. If they did this, their property theory would not be based on labour, like Daves or Locke's; but on conflict-avoidance, like the modern libertarian theory.
\section{The Mises Caucus as a Cargo Cult}
\label{sec:org0e33688}
The Mises Caucus have a lot to say with respect to their recent pragmatism:
\begin{quote}
I think sometimes libertarians struggle to adjust to the facts on the ground, [\ldots{}] I think a lot of people in the libertarian world suffer from ideological possession, it's almost like: dealing with the changes in the real world, that's off to the side, we're talking in theory [\ldots{}]

All I'm saying is that Donald Trump is something different.\footnote{TomWoodsTV, \emph{Dave Smith \& Tom Woods on Voting Trump | TWS \#2563}, t. 12:17}
\end{quote}

Trump is something different, and that is the end of the discussion. Again, what works today might not work tomorrow! We must be in a perpetual state of try and see. The standard of the true is what works to satisfy the arbitrary whims of a consciousness. Reality is controlled by thought.

They like to claim that this has always been the Mises Caucus position:
\begin{quote}
I do understand where the strategy that Angela kind of embarked on here in the libertarian party, and a lot of us libertarians [\ldots{}] this was very much a different strategy for all of us. [\ldots{}] I would just kind of implore a lot of libertarians to think about the fact that the world has changed. Truly drastically the world has changed over the last few years. [\ldots{}] the world has changed, and in a situation like that you should at least be open to new strategies. Rather than having this, what seems to me, quasi-religious allergy to even suggesting that maybe just sitting in our corner and preaching our message is not the best use of this party.\footnote{Dave Smith, \emph{Did I sell out my principles by supporting Trump?} (Twitter Space), \url{https://x.com/ComicDaveSmith/status/1860865805006590365}, t. 2:30}
\end{quote}

So supporting Trump is ``something new.'' The Libertarian Party has really been controlled by radicals this entire time, and now we need to set out on a new, pragmatist path and see if that works. Again, the truth is just what works, it is entirely malleable depending upon the ``facts on the ground'' for these people. It doesn't matter that the LP was literally under the control of the prag caucus for decades and that the Mises Caucus was a revolution \emph{against} pragmatism. That was then and this is now; so now its the case that the LP used to be radical and we now need to try out pragmatism for a bit.

They want us to sincerely believe that they are just following the greats:
\begin{quote}
You know, Hoppe's strategy, if you actually read the entire piece, has been tremendously influential on Michael Heise's, uh, decentralized rebellion, right? [\ldots{}]

This is the thing where we have to marry our principles to the real world. And I will tell you a lot of these kind of like, and believe me, I've said them all in my day, but a lot of these anarchist mantras of we just have to pull out. We can't work within the system. We can, it just becomes an excuse to do nothing and say that like, oh, yeah, so we don't actually, I don't actually have to go grind it out and do the hard work to maybe make things 1\% better.

I can just retreat and say, oh, yeah, you're, you know, painted now, you're not as pure as me. This is what having a conversation about strategy, strategy comes down to the real world, what are plausible, realistic outcomes, and what outcomes are preferable to other plausible, realistic outcomes. That's what we're talking about here. Ross being freed is better than Ross not being freed, and it is a very big deal to win.\footnote{Dave Smith, \emph{Did I sell out my principles by supporting Trump?} (Twitter Space), \url{https://x.com/ComicDaveSmith/status/1860865805006590365}, t. 1:45:38}
\end{quote}

So, let's look into the strategy forwarded by our libertarian greats, and see whether it lines up with the pragmatism now being espoused. Murray Rothbard is surely the greatest of the greats in this regard. He is the man behind the ``Leninist'' strategy described above.

As the neo-prags will gleefully point out, Rothbard did indeed build coalitions with rightists toward the end of his life. They take this fact to conclude that we must also ally with rightists, whatever the cost. Like the cargo cults attempting to emulate some distant effect of air travel by building their wooden airports, the neo-prags attempt to imitate the Rothbardian strategy by engaging in the distant effect (coalition building) without understanding the more fundamental strategic causes that imply such coalition building to be proper.

What these people fail to understand is that Rothbard always kept principles at the core. He would vociferously refute any inkling of pragmatism the moment it reared it's ugly head. Certainly, he would build \emph{coalitions} with people, but never to the detriment of the ideal. We could power a mid-sized city with the torque that must be being exerted by Rothbard rolling in his grave over these imbeciles using his name to advocate fundamentally anti-libertarian policies for the sake of ``practicality'':
\begin{quote}
Since anarchists and other libertarians are, to say the least, an embattled minority, we have tended to be indulgent toward anyone and everyone in our ranks, even those who have been busily pecking away at the vitals of the libertarian position; Or, to change our metaphor, who have been picking off stragglers as we try to continue our march against the State. [\ldots{}] I say it is high time to take off the gloves in the struggle against those whom Sam Konkin incisively calls ``Natural Outlaws.''

For many years now, first as anarcho-Stirnerites and now also as anarcho-pragmatists, the Natural Outlaws have been firing away from inside their supposedly impregnable fortress of ethical nihilism, sneering at such fundamentals of libertarianism as individual rights and rational ethical principles. Since the Outlaws stand for nothing, they believe they can remain permanently blessed with the advantage of the strategic offensive.\footnote{Murray Rothbard, \emph{On The Duty Of Natural Outlaws To Shut Up}, \url{https://mises.org/articles-interest/duty-natural-outlaws-shut}}
\end{quote}

Rothbard would go on to outline a line of attack against these natural outlaws in order ``that ethicists and natural lawyers [may] take the strategic and tactical offensive.''\footnote{Murray Rothbard, \emph{On The Duty Of Natural Outlaws To Shut Up}, \url{https://mises.org/articles-interest/duty-natural-outlaws-shut}} Namely, he shows that \emph{on their own premises}, the pragmatists have a moral duty to shut up and abandon their moronic creed:
\begin{quote}
First, the trouble with pragmatism, and especially anarcho-pragmatism, is \emph{that it doesn't work}. And since pragmatists believe that the only truth is whatever ``works,'' that settles that (setting aside such deep problems as the meaning of ``work:'' work for ``what,'' etc.). Take, for example, the severe criticisms that Jorge Amador, the guru of anarcho-pragmatism, has made of the Bergland-Lewis presidential campaign, in his organ, The Pragmatist. Amador's critique is that the LP was (a) too gradualist, and also (b) ideological. In other words, his preferred campaign would be radical anarchist to the hilt, and yet non-ideological. That is, all talk of moral principles or rights would be tossed aside. Not only that: we could no longer call the State an organization of a criminal ruling class, because ``crime'' itself is a moral and natural-law concept, and presumes immoral criminals ripping off innocent victims. So what would an Amadorean LP campaigner talk about? He would confine himself to demonstrating the pragmatic virtues of the radical anarchist alternative.

But this is a tall order indeed. In fact, a virtually impossible one. The pragmatic radical anarchist is faced immediately with powerful critiques from pragmatic statists. He can show, for example, that anarchy would increase production, yield a higher standard of living, etc., in the long run. But in the short run, lots of the privileged, subsidized, or monopolistic would be cast adrift. All these short-run and maybe intermediate-run problems could only be offset by vague future benefits. But why pragmatically, should everyone prefer the long-run to the short-run? What about the high-time preference people, who thus challenge the Amadorean: ``Look here, fella, I \emph{know} the pragmatic benefits I’m getting from the current system. And I \emph{know}, too, the headaches, the disruptions, the losses that I and lots of others will suffer during the lengthy `transition' period. Even if you've convinced me that eventually I might benefit, these benefits are too chancy and too long-run for me to want to risk it.'' And if the average person cannot be sold on radical immediatist anarchism, \emph{a fortiori} the criminal ruling class, those net beneficiaries of the State, they who might well be losers even in the long-run, \emph{certainly} won't be convinced. \emph{At best}, the Amadorsymp will say: ``Well, I admit this anarchism sounds pretty good. But pragmatically, to ease the transition and minimize the costs that even you admit, let's move toward the ideal very, very gradually.'' And we are back, will-nilly, to the Republication or Democrat Party, the master ``gradualists'' of us all.

It is no accident, then, that Democrats and Republicans proudly \emph{call themselves} ``pragmatists.'' Sure, they believe in freedom, in peace, in free markets, in all the goodies, but these goals have to be approached, they tell us, piecemeal, by the groping push-and-pull of the democratic consensus. And we are back hip deep in the status quo. ``Radical pragmaticism'' of any sort, whether anarcho or Khomeini or whatever, is virtually a contradiction in terms.\footnote{Murray Rothbard, \emph{On The Duty Of Natural Outlaws To Shut Up}, \url{https://mises.org/articles-interest/duty-natural-outlaws-shut}}
\end{quote}

It is entirely expected that these high time preference pragmatic arguments would arise; it is demanded by the very philosophy of pragmatism. If what works today might not work tomorrow, there is a philosophical imperative to prefer today over tomorrow. We have no idea if by the time that the long-run benefits of anarchism are supposed to kick in that anarchism will still even ``work.'' Heck, given the standard of ``working'' is that it satisfies the purposes of a consciousness; there are certainly pragmatists who find themselves in a position where they wish to violate the rights of others by sending police to break up drug deals, or sending the entire US military to prevent immigration. These pragmatists would be sworn, by virtue of their aggressive whims, to reject anarchism because it could not work to achieve what they desire. The correct response is, again, to reject the pragmatist basis that the standard of the true is that which works to achieve some end that you arbitrarily hold.

It simply cannot be stated enough times. \emph{Pragmatism is the opposite of principles.} There is no middle ground, or dialectic synthesis, or compromise possible between these two extremes. One can either adopt the radicalist strategy advanced by Murray Rothbard and be bound to immutable ethical principles; or one can abandon principle and truth, descend into pragmatism, and in no way be properly described as a libertarian.

See for yourself the effects that this neo-pragmatism of the Mises Caucus has had. They used to make fun of the LP for trying to cozy up to power to get a seat at the table; they talked about how the only reason Cato types get to go to cocktail parties with people in power is because they don't consider them to be an actual threat. Now history is repeating itself, but with the right instead of the left. \emph{Both} these beltway Washington libertarians and the neo-prags enable, rather than meaningfully oppose the regime. All pragmatists necessarily enable the philosophical status-quo.

Dave Smith's right-pragmatism has had the effect of placing libertarianism even lower in the food chain than the paleocons; who are themselves already fighting for scraps as the neocons get whatever they want. Dave's job becomes trying to pander to paleos and other dissident rightists and try to convince them that their end goals are easier to achieve without the state. If he were to talk with Ibram X Kendi, there is no way he would approach it by saying ``look, I know you want to implement DEI policies, but in a libertarian society you can achieve these things with covenant communities!'' This is because it would be very easy for him to recognise that these progressives are completely disinterested in our philosophy and have absolutely no intention of living and letting live. The Ibram X Kendi's of the world would not be willing to just leave white people alone, they want to actively oppress white people because of their Marxian race theory, which places revolution against the white man as a necessary next step in the evolution of mankind.

Tom Woods cogently points out the problem with this sort of pragmatism when asked about the removal of the ``bigotry is irrational and repugnant'' line from the LP platform:
\begin{quote}
I suspect that what they're driving at here is that this is more of an attempt to sort of curry favour with people who were never going to like us anyway.\footnote{ReasonTV, \emph{Inside the Mises Caucus Takeover of the Libertarian Party}, \url{https://www.youtube.com/watch?v=NsgFdPqOAhk}}
\end{quote}

Tom is entirely correct about this, but this applies both to rightists \emph{and} leftists. The neo-prags are able to apply this standard only when it comes to leftists, whenever Dave is faced with any sort of dissident right-winger, he just can't help but to try and frame our philosophy as a solution tailored exactly to their problem-statement. Instead of attacking the root problems that exist in, say, fascism, the Dave Smith school of debating rightists requires that you accept their whole framing and try desperately to please them.
\section{The Border Question}
\label{sec:org7b0ebaa}
\subsection{The Neo-Prag Case for Immigration Restrictions}
\label{sec:org0273d9c}
The main omission up to this point in the video, which I have hinted at several times is the big reason for this pragmatist shift. It is no coincidence that neo-pragmatism was able to sweep the de-facto leadership of the libertarian movement in such a short amount of time. These people were all infected by this culture warrior brain worm because of their stance on the question of immigration controls. Namely, they want the state to restrict immigration, and have worked themselves into knots trying to justify this stance.

I shall take as the definitive celebritarian argument for closed borders to be the one that Dave presented in his debate with Alex Nowrasteh. Of course, the neo-prags take great inspiration from Hoppe, but I regard the neo-prag hordes as being far too weak on theory to have any chance of using his arguments at any level beyond that of a Hoppean Cargo Cult. In any case, I have already addressed Hoppe's argument in my moral case for open borders;\footnote{LiquidZulu, \emph{The Moral Case for Open Borders}, \url{https://www.youtube.com/watch?v=I3sa7J-i8ZY}} so for now I will allow Dave to make the neo-prag case. I choose this debate in particular, because first it is relatively very recent, so presumably it is the most up-to-date version of the argument; and Dave had months to prepare this argument, rather than having to come up with it off-the-cuff, which we all know he is so good at.

To be clear, the proposition that Dave will be attempting to prove here is that ``government restrictions on the immigration of peaceful and healthy people make sense from a libertarian standpoint, especially in present-day America'':
\begin{quote}
Really the foundational principle of libertarianism is [\ldots{}] voluntary exchanges. [\ldots{}] Even things like property rights, the NAP, all of them are built around the idea that we prefer voluntary exchange, and we reject aggression.\footnote{ReasonTV, \emph{Dave Smith and Alex Nowrasteh debate immigration}, \url{https://www.youtube.com/watch?v=6xWiSBnwIQM}, t. 3:07}
\end{quote}

So, right off the bat, this is an absolutely abysmal starting point; perhaps even worse than his previous ``retroactive slavery'' theory of property. He says that the entire foundation of libertarianism is ``the idea that we prefer voluntary exchange, and [\ldots{}] reject [aggression].'' He says that property rights and the NAP are not essential themselves, they are rather derived from this preference for voluntary exchange. Of course, this is entirely consistent with his pragmatist philosophy. He can't rely on the correspondence theory of truth to yield some immutable ethical principles or anything like that. The basis of validation itself is some arbitrary whim of his. So accordingly his legal theory is rooted in him merely \emph{preferring} voluntary exchange to aggression. At least he has started being honest about this and baking it directly into his property theory.

At any rate, let's see how he will apply this preference-theory of property to the case of the border:
\begin{quote}
If that's your starting point and then your end point is supporting a policy [that would] drastically change almost every aspect of our society against the will of almost all of the people who live here. I would humbly suggest that you have messed up somewhere along the line.\footnote{ReasonTV, \emph{Dave Smith and Alex Nowrasteh debate immigration}, \url{https://www.youtube.com/watch?v=6xWiSBnwIQM}, t. 3:59}
\end{quote}

Indeed; if just base whim is the standard of truth and justice, and Dave just prefers voluntary exchange to aggression; then we probably have messed up if we are to conclude that Dave would prefer allowing people into the country. But libertarianism is not based in the primacy of Dave's consciousness. If you have a rational theory of property based in conflict avoidance, you can quite readily conclude that the state doesn't have the right to exclude people from the United States of America. The state doesn't own the United States of America, and accordingly, you initiate zero conflicts by immigrating to the United States of America. At no point in such a proof do you have to consult ``the will of the people'' who live in ``our country.'' Countries are not a valid legal category, the state is a gang of criminals; and there is no collective will of the people that one may or may not consult.

So it's already out on those two fronts, but even if we are to attack this through the ``voluntary exchange'' lens, there is still no problem with open borders. That a company wishes to hire Arjun from India \emph{is a voluntary exchange}. The company is voluntarily giving him a wage, and he is voluntarily working for that wage. Literally every other person in the US could be upset about this and it would not matter. The \emph{only} two people who matter are the people who are actually doing the exchange. Certainly, if every single person in the United States didn't want any immigrants to come over, having open borders wouldn't mean shit! If I were to try and immigrate in such a world, who would sell me a house? Who would sell me food? Who would hire me? Americans are clearly \emph{already} voting with their dollars to bring in enormous numbers of immigrants. To be sure, state interventions absolutely warp this and cause forced integration; but the solution to forced integration \emph{is not} forced \emph{disintegration}.

That is, it is true that such policies as welfare and trade restrictions\footnote{See: Hans-Hermann Hoppe, ``The Case for Free Trade and Restricted immigration,'' \emph{Journal of Libertarian Studies} 13:2 (Summer 1998): 221-233} (such as tariffs) do indeed increase the incentive to immigrate; and it is further true that civil rights legislation prevents people from peacefully disassociating from the horde of immigrants that are sure to come. But that does not imply that we, as libertarians, should focus on trying to get a government solution to what is decidedly a government problem. We should use the public stage that the immigration issue has predictably achieved to tell people that we need to shrink the welfare state, remove barriers to trade, and do away with the disastrous civil rights regime. The weakling libertarians like Dave failed this easy test, and instead used the opportunity to sell out libertarianism to the very people who are exacerbating the issue.

In any case, Dave continues:
\begin{quote}
If you look at where the American people are [\ldots{}] almost no pollster will ask: ``do you support open borders?'' [\ldots{}] But if you look at where the American people are on immigration, it's nowhere near there. And it's about as much of a consensus on any issue as you could find in America that people reject Alex's position on this. That doesn't mean that Alex is wrong, [\ldots{}] I'm not making an appeal to democracy or an appeal to popularity, I'm making an appeal to voluntarism. My argument is that if the people of a country do not wish to have a large number of immigrants coming in, then it is not a voluntary transaction to force that on them.\footnote{ReasonTV, \emph{Dave Smith and Alex Nowrasteh debate immigration}, \url{https://www.youtube.com/watch?v=6xWiSBnwIQM}, t. 5:00}
\end{quote}

Ok so, Dave's argument is that the ``American people'' do not support open borders, and this means that it is not a voluntary transaction to ``force'' this on them; and this is supposed to have been established \emph{without} any appeals to popularity. But who exactly is having immigration ``forced'' on them? Dave concluded that the ``American people'' are having immigrants ``forced'' on them because the ``American people'' don't like immigrants. That is, immigration is unpopular. Certainly the individuals who are actually hiring immigrants aren't having it forced on them, so what is he left with without the popularity argument? This is a lot like if I were to say that I in no way wish to make an argumentum ad hominem against Dave's position here, but that being said Dave is incorrect because he is a comedian rather than a philosopher. Now my prefacing of the claim with saying that it isn't an ad hominem attack might be enough to confuse the gaggle of morons whom Dave wishes to court; but it doesn't actually make it not an ad hominem attack.

Nowhere else in libertarian theory do we make such appeals to popularity. The vast majority of renters and home owners believe there should be rent caps,\footnote{\url{https://x.com/cremieuxrecueil/status/1943433060609069414}} but that does not establish that rent increases are being ``forced'' on renters. Most Americans are opposed to selling off public land to private developers, but that does not establish that they are having these sales ``forced'' on them---on the contrary, the government giving up control of these things is it \emph{releasing} it's force in that area. This is why ``force'' is simply a sloppy and vague term when it comes to legal analysis, and why it should never be used. If people voluntarily hiring immigrants, selling to immigrants, becoming friends with immigrants counts as ``forcing'' immigration on the nebulous ``American people,'' then ``force'' is clearly not the correct concept to use. With the proper conflict-avoidance theory of property we can see the truth that in no way is immigration an aggression against anybody.

Furthermore, how can Dave apply this collective will analysis to the blue strongholds that are overwhelmingly opposed to ICE? Does Dave want the federal government to force immigration restrictions on sanctuary cities where the collective will is to allow immigration? Are the massive anti-ICE protests and entire apps whose purpose is to track the movements of ICE agents not ample evidence that there isn't some unified American will? Does this not destroy Dave's entire collectivist claim?

Regardless, Dave presses forth:
\begin{quote}
It certainly does seem, as a reasonable starting point, [\ldots{}] that we don't want the government doing anything to anyone. You know, we don't want the government stopping someone who's not aggressing, [\ldots{}] but the truth is that we find ourselves in an awkward position when it comes to immigration. Because immigration is a government policy whether you have open borders, or closed borders. [\ldots{}] As long as the government is controlling those borders and the government is dictating immigration policy what you have is a government program.\footnote{ReasonTV, \emph{Dave Smith and Alex Nowrasteh debate immigration}, \url{https://www.youtube.com/watch?v=6xWiSBnwIQM}, t. 6:20}
\end{quote}

This is a point that I see basically exclusively being used in the immigration debate. I don't think that for a second Dave would permit the argument that ``well, so long as the FDA exists, any policy on drugs is still a government program, so given that most people want fentanyl banned it should be banned.'' This is a sloppy equivocation on the word ``policy.'' Quite simply: not all policies are actually in the same category. Consider the 2nd Amendment. This is a ``policy'' that says the government is not allowed to make any policies as far as the right to bear arms is concerned. This is not the same as ``policy'' in the ordinary sense of the word. This is a ``null-policy'' that is a rule that says the government is not allowed to make any rules in some area. As libertarians, we must advocate null-policies as a matter of principle.

That null-policies are grouped together in a big package deal with actual positive policies is nothing more than a result of legal positivism. If all law comes from the government then its all of the same ethical significance: there is no natural law that the government can become more aligned with by negating policy. Open borders is the null-policy here, it is not ``just the same'' as restricting immigration. The open borders position is that the government should not be allowed to make any rules pertaining to immigration. That is not the same as saying that the government should make certain rules pertaining to immigration. This is not difficult, Dave.

Regardless, Dave goes on to point at exactly the problem I have with him only ever debating people who are not radicals on the border issue:
\begin{quote}
It is impossible to oppose this resolution, unless you support unrestricted open borders, period. [\ldots{}] Because even if you are going to argue [\ldots{}] for an ``Ellis Island'' system, you're still supporting an enormous restriction on peaceful and healthy people. [\ldots{}] It's almost like they're kind of admitting that immigration is different. It's a different situation. [\ldots{}] So often people who are for open borders won't bite the bullet and say there should be absolutely no restrictions, but they argue as if they do believe that.\footnote{ReasonTV, \emph{Dave Smith and Alex Nowrasteh debate immigration}, \url{https://www.youtube.com/watch?v=6xWiSBnwIQM}, t. 8:40}
\end{quote}

Indeed; this leaves the various ``GDP go up'' type of arguments you see coming out of the CATO institute and other such organisations as being vastly weaker than a true ethical defense of open borders. And this certainly makes one wonder why Dave utterly refuses to debate me, when I am surely exactly what he is looking for here. I am perfectly willing to bite that bullet along with many stronger ones. Perhaps he could learn a thing or two and see how its done; considering his own allergy to bullet-biting.

\begin{quote}
This all comes down to the issue of government property. [\ldots{}] Many libertarians will act as if its a given that there should be no restrictions on government property. But that simply isn't true, it's not deduced from libertarian principles and none of us would apply that standard for life in general.\footnote{ReasonTV, \emph{Dave Smith and Alex Nowrasteh debate immigration}, \url{https://www.youtube.com/watch?v=6xWiSBnwIQM}, t. 11:10}
\end{quote}

Of course, Dave is just purely wagering that this stance isn't deduced from libertarian principles, because he doesn't actually know what the libertarian principles are. So let me help him out and construct a purely deductive proof that there should be no restrictions on government property.

First, for any scarce means in existence, there are three possibilities: (1) nobody owns it; (2) one person owns it, or; (3) a group of people own it. If nobody owns it, then nobody has the right to exclude anyone from it, including immigrants. If one person owns it, then that one person is allowed to exclude anybody he wants from it. And our remaining possibility, that a group owns it implies a contradiction,\footnote{See: LiquidZulu, ``On the Impossibility of Group Ownership,'' in id., ``Homesteading and Property Rights,'' in id., \emph{The Fundamentals of Libertarian Ethics}, \url{https://liquidzulu.github.io/homesteading-and-property-rights}} so it is out. To demonstrate, take the example of a group of people, \{A,\ldots{},Z\} each commonly owning a stick. A decides he wishes to use the stick to spearfish, but B does not want this to take place. So, there is a conflict over the use of the stick, and who should win this conflict? Well, A owns the stick, so A should win the conflict, and the spearfishing should go forth. Also, B owns the stick, so B should win the conflict, and the spearfishing should not go forth. So the presumption that both A \emph{and} B own the stick implies that the spearfishing both should and should not go forth; a contradiction.

So the only possibilities for some government possession are that a single person owns it, or that nobody owns it. Obviously in cases of eminent domain where the government literally steals some property from a single, identifiable person; then yeah, that person should be allowed to tell any immigrants to stay the heck away. But that is not the case for the vast majority of public possessions. Indeed, the burden lies on Dave to prove that for some particular possession it is owned by some one particular guy who doesn't want some particular immigrant to use it. The presumption of innocence comes from the basic epistemic principle that the burden of proof lies on the one making the positive claim. Without such proof, the assumption is that the immigrant is allowed to go there. So it has been established purely based on the libertarian conflict-avoidance theory of property that the presumption is that there should not be any restrictions placed on government property.

Dave attempts to attack this stance with a series of reductios ad absurdum:
\begin{quote}
My point is that many of the people who argue for this lax immigration policy they act as if that's the pure libertarian position, but ask yourself: is that really the pure libertarian position? Is it the case that if the government is currently in control of some property that therefore there ought be zero restrictions? If that were the case, then public schools couldn't tell a random 55 year old drug addict that he can't enter the school, and hang out in the girls bathroom. [\ldots{}] Would you say that you should be able to scream obscenities in a public library, that you should be able to lay down and camp out in the middle of a freeway? Is the libertarian position that the government should not be allowed to implement a speed limit? I don't think that actually follows.\footnote{ReasonTV, \emph{Dave Smith and Alex Nowrasteh debate immigration}, \url{https://www.youtube.com/watch?v=6xWiSBnwIQM}, t. 12:30}
\end{quote}

Now, with these reductios Dave is employing a very sneaky trick. Namely, if there truly were \emph{no} restrictions on government property, then it wouldn't be just trucking along like normal except with more hobos. Very quickly these state possessions would be homesteaded and allocated to their most value-productive use. Dave isn't actually talking about \emph{no} restrictions, he is trying to sneak a series of restrictions through the sally port. Namely that the government is preventing homesteading of it, forcing people to use it, and thus preventing them from excluding unsavoury characters. Indeed, it would be bad if the government continued to forestall these things, but the open borders position is precisely that the government \emph{not} forestall them any longer.

So, indeed, if the most value-productive use of school bathrooms is to be a hang-out spot for homeless drug addicts, then great, that's what they will be used for. If libraries would be better used as obscenity-screaming zones, then entrepreneurs would find this out. Dave, on some level, believes in the central planning that has yielded the current uses for these items. He thinks there is some metaphysical teleology that Big Brother is able to use to divine that what is being used as a school should stay as a school; and what is being used as a road should stay as a road; and so on.

Indeed, it has been pointed out in this debate that border control is an utterly fundamental aspect of the state. I completely agree, and that that is the case is revealed by \emph{actually} making the reductio scenario of no restrictions on government property. If you truly imagine a state without \emph{any} border restrictions \emph{at all}, you quickly find that the state pretty much evaporates. Sure, it might remain as a nomadic set of disconnected barbarian tribes with no permanent base of operations; but this is nothing like what we understand the state to be. Of course, the bordertarians try to use the point that borders are so fundamental to the state to say that we need to put off abolishing them until we deal with other things like welfare. But is the opposite not the case?

If the state was eviscerated down to the level of tiny nomadic gangs, would they in any way be able to maintain such a thing as a welfare state, or a standing military, or a federal tax collection apparatus? Would we apply that same argument to any other aspect of the government? Would we say that because the income tax is so fundamental to the operation of the US federal government, we should put off abolishing the IRS until we get other things sorted out? No! We would press that button until our fingers had blistered. It's just that on the border issue, the liberty movement has almost universally fallen for the government propaganda---that without them keeping ``our border'' secure terrible things would ensue.

These Schrödinger's bums, who are simultaneously so lazy that they can't hold down any employment and rely on the welfare state to stay alive; but also upon the removal of government restrictions would mass mobilise to take over  Central Park, simply do not exist. Even if those bums did pull off such a feat, would private investors not immediately buy it off of them to turn into condos or whatever? Why do the Dave Smith's of the world completely forget the basic economic law that upon the introduction of market forces, \emph{any} scarce good will tend towards its most efficient allocation?

\begin{quote}
I think that if we recognise that taxation is theft, and we recognise that governments are built off of aggression: they built and maintained public property with tax dollars that they took at the threat of violence from the American people. And what the American people want them to do with those tax dollars ought to play somewhat of a role. Now I'm not arguing as some other libertarians have, and I think they overstate the case a little bit, that public property is ``owned'' by the taxpayers, or ``owned'' by the net-taxpayers. I'm simply arguing that the American people have a better claim than others do. Which is really, I think, what property rights starts as. The idea that if the American government is doing something it ought be for the American people, it ought be in their interests. If there is a local public school and there's a group of people in the community who have been paying property taxes into the school for many many years, I think that school ought serve those people's interests, and not equally the interests of the world.\footnote{ReasonTV, \emph{Dave Smith and Alex Nowrasteh debate immigration}, \url{https://www.youtube.com/watch?v=6xWiSBnwIQM}, t. 13:28}
\end{quote}

OK, so the ``American people'' ought to have ``some say'' over how public property is managed; but they don't ``own'' the public property. Then, \emph{what} say should they have over it, Dave? What specifically should ``they'' be allowed to do with it? What does it mean for a collective like ``the American people'' to even have an opinion at all, let alone an opinion or a say in how some particular good is being utilised? Even if some particular American taxpayer has a greater claim to a piece of government property than a random person on the other side of the globe, this does not establish what Dave needs it to establish. All that would establish is that in a conflict between that particular taxpayer and that particular random other person over the use of the government property; the taxpayer should get precedence when it comes to homesteading and management. This does not establish that the government should keep centrally planning it in a way that they think might align with some non-existent collective will.

Indeed, the way for that local school to best serve the interests of the community in its proximity is to have the government no longer enforce any restrictions on its use. If the government stopped forestalling it, immediately it would start to be allocated towards its most value-productive use; and the local community would be able to extract the maximum \emph{just} benefit from it. Indeed, some members of the local community might lose out; maybe they want to keep being a welfare leech and educate their kids for free, or they want to keep their cushy government job as a molestor of the minds of children. That those people would lose out isn't relevant---it is good and proper that such monsters lose out, and it is time for libertarians to embrace this fact. We want every single public school teacher to be fired, and have all of their assets seized to compensate the taxpayer. It is just a shame that no method of brain damage exists to properly retribute for the countless children they have mutilated the minds of during their career.

The losers of such a mass privatisation campaign as open borders points at the true motivation behind these immigration restrictions and why they are so popular among the native population of pretty much everywhere. Namely, immigration restrictions are a form of protectionism. Natives often don't want immigrants to come in and compete for jobs with them; so they turn to daddy government and ask that he keep out this competition. This is why in due course, the podcast bros whom Dave told us he was swaying to liberty, have since run back to Bernie Sanders and Zohran Mamdani.\footnote{\url{https://x.com/TheTNHoller/status/1943433346790412296}}

This should not be surprising to anybody who has paid attention. The ``based and redpilled'' rightists are just as socialist as the leftists, they only focus on different issues. People like Dave seem to only be able to recognise this when it comes to war; but whenever the socialism extends to his pet issue of immigration, he is blind as a bat.

In any case, he continues:
\begin{quote}
There is nowhere in libertarian principles where you can deduce that because I stole something from you, it is now unowned, or owned equally by everyone in the world. If taxation is theft, and everything the American government does has been achieved by robbing the American people, then the rightful owners are still the American people. Now, I'm not saying it's completely possible to divvy that up perfectly, but you know what's definitely possible? To not do what 95\%+ of the American people don't want which is to open the borders.\footnote{ReasonTV, \emph{Dave Smith and Alex Nowrasteh debate immigration}, \url{https://www.youtube.com/watch?v=6xWiSBnwIQM}, t. 14:38}
\end{quote}

Indeed, the government stealing some tax money from a person doesn't imply that that money is now unowned, or owned equally by everyone in the world. I will be charitible and assume that Dave has seen this argument being made by someone at some point; but that certainly isn't my argument or the argument of any other radical open borders proponent. Indeed, to the extent that you were robbed by the state, you still own that money. But what you \emph{don't} own is some local park that the government later spent money on---certainly not in tandem with every other taxpayer.

It is on this point that Dave's stance is entirely confusing. He says on the one hand that the taxpayers don't \emph{own} government property as a group, but then what is it that he thinks is still owned by individual taxpayers if not the \emph{money} that was taken from them? There simply is not a valid chain of title-transfer that would yield anything else. When the government stole that money and paid some contractor to build a park; the money is still owned by each individual it was stolen from. That trade was not a valid transfer of the title to property.

And heck, let's take Dave's standard a little further here. He says that given you can't perfectly divvy up everything that has been stolen, therefore what should be done is to not have open borders because 95\%+ of American people don't want it. Is that a valid standard anywhere else? Given the fact that the vast majority of Americans want some gun restrictions, especially in urban centres; should the government restrict the ownership of firearms on any public property? Should I be stopped from importing a firearm to Manhattan because the vast majority of Manhattan residents don't want me to? That doesn't strike me as libertarian at all.
\subsection{Immigration Restriction Leads to Tyranny}
\label{sec:org7b4ab46}
So much for Dave's presentation of immigration restrictionism, then; allow me to explicate what exactly his policy proposal practically results in. Certainly, Dave's proposal would require massive government force to implement. He recognises this openly:
\begin{quote}
We've gotten to a point where it's like: yo this is fucking insane! Like, I don't care; if you believe in open borders right now under current circumstances, you're an insane person and you're as bad as a communist. You're as bad as saying: well, I believe in this theory, and sure my theory is going to result in the destruction your kid's lives, but I'm so committed to this theory that you've gotta break a few eggs to make an omelette. I'm just not even interested in talking to you anymore. [\ldots{}]

The answer is: take our fucking entire military, and put them on our border, and secure our border. [\ldots{}] For all the shit the US government is capable of doing--and it is the biggest, largest, most powerful organisation in the history of the world--[\ldots{}] if you think the strip from fucking Texas, Arizona, and California, over Mexico can't be shut off, like we don't have the power to do that, I'll tell you, we absolutely do, we would just need the political will to do it.\footnote{Liberty Lockdown, \emph{LIVE REACTION with Dave Smith: Did the controlled narrative just collapse?}, \url{https://www.youtube.com/watch?v=Erm59oQn--8}}
\end{quote}

Dave, the issue has never been about the \emph{ability} to do it. Certainly, the US military would be able to use all of their tanks, and guns, and bombs to raid every home in the 100-mile border exclusion zone and declare martial law to completely shut down any chance of immigration. That's not the issue with it. The issue is that this would be completely abhorrent by libertarian standards.

And martial law is not an exaggeration, here. Let's consider the actual practical implication of setting the \emph{entire US military} at the task of enforcing immigration laws. Let's ``live in reality'' as Dave constantly asks us to. This policy would involve military checks of every business in the US on a regular basis---why else do you think that everyone has to sign a bunch of forms before they can get a job? It's so the business can prove that it is employing only people who have government permission to work, rather than illegal immigrants. It would involve military checkpoints on pretty much every major road they could get their grubby mitts on. It would involve military raids of residential properties; predator drones flying overhead; tanks rolling down 5th Avenue. This would be the single greatest rise in militarisation in US history. If you thought Tiananmen Square was bad, how about copy and paste that to every single major city in the US; most of whom have a primarily Democratic population who would be opposed to these restrictions.

This is an especially interesting take for Dave to give in a podcast, the entire topic of which was how bad the American war machine is. How the British and American militaries being placed on the Iraqi border was a prelude to war with Saddam---a fact that was identified \emph{in this very podcast}. What does the great Dave Smith, most consistent motherfucker you know, think that placing \emph{the entire} US military on the Mexican border might lead to? It's not like Trump declared cartels as terrorist organisations and suggested invading Mexico to take them out \emph{already} or anything.

Simply take Dave's exact same proposal here and apply it to a different country, Israel. Dave has spent a great deal of his career as of late shouting from the rooftops over how terribly the Palestinians are being treated by the IDF. But, given his policy for America, should he not be in favour of having the entire Israeli military be set at the task of enforcing the border claims of the Israeli government? Is this not the exact policy that Israel is doing which has led to all of the terrible things he is making a career pointing out? Should we take a military many times larger than the IDF and have it do the exact same things as the IDF but in the US instead? Does Dave Smith trust them to do that without ending up exactly the same as the IDF, just way bigger?

Notice also the framing he gives us: radical anarcho-capitalists are ``as bad as communists'' because both the radical anarcho-capitalists and the communists are just too attached to theory. This is a total equivocation. The whole reason why the communists come to so many terrible, evil conclusions in ethics is \emph{because they reject the correspondence theory of truth}. To paint the defenders of that theory of truth with the same brush as communists \emph{whilst also} rejecting the correspondence theory of truth in favour of the pragmatic one is utterly sickening.

This is not merely an off-the-cuff, poorly worded rant by Dave either. To be sure, it is that; but not \emph{only} that. Dave would double down on his sentiments here a few weeks later:
\begin{quote}
I'm really glad that I said what I said on Clint's podcast\footnote{The Biblical Anarchy Podcast, \emph{What is the Libertarian Response to Immigration w/Dave Smith}, \url{https://www.youtube.com/watch?v=ghqN\_F2iTdU}, t. 07:34}

{[}\ldots{}]

I don't take it back to me that's as bad---and look; to be clear, when I said to that Pho lady that you're as bad as a commie or a nazi, [\ldots{}] I mean you're as bad as someone in America today who is a nazi or a commie.\footnote{The Biblical Anarchy Podcast, \emph{What is the Libertarian Response to Immigration w/Dave Smith}, \url{https://www.youtube.com/watch?v=ghqN\_F2iTdU}, t. 21:00}

{[}\ldots{}]

Calling you a commie or a nazi is the nicest thing I can say. What I actually would say, I feel uncomfortable saying on air. [\ldots{}] I'll just say this, and I mean this, there are some things worth killing people over, and there are some things worth dying over, there are some things worth dying in order to have a chance to kill someone over, and for me that's it.\footnote{The Biblical Anarchy Podcast, \emph{What is the Libertarian Response to Immigration w/Dave Smith}, \url{https://www.youtube.com/watch?v=ghqN\_F2iTdU}, t. 22:37}

{[}\ldots{}]

There were a lot of people who got really upset about the way that I said it on Clint's podcast, and I will admit it was a night-time podcast, where I was having a whiskey [\ldots{}]. I stand by everything I said. But when I said ``bring the entire military home and put it on the border,'' sorry, that's a fucking win.\footnote{The Biblical Anarchy Podcast, \emph{What is the Libertarian Response to Immigration w/Dave Smith}, \url{https://www.youtube.com/watch?v=ghqN\_F2iTdU}, t. 1:11:18}
\end{quote}

He would re-iterate this again \emph{after} the totalitarian nature of Trump's immigration enforcement became clear to everyone:
\begin{quote}
Andrew Wilson: I have to say, the highlight of my evenings has been watching the cops beat the fuck out of these people every single night, and it's with great joy that I watch this, does that make me a bad person, Dave?

Dave Smith: No. I'm all for shutting this shit down by any means necessary.\footnote{TheCrucible, \emph{The Extravaganza EP: 03 with Dave Smith | Feminist Flips Out | Riot Bloopers | Real Jews Revealed (6/11/25)}, \url{https://rumble.com/v6um393-the-extravaganza-with-special-guest-dave-smith-61125.html}, t. 15:00}

{[}\ldots{}]

The mass-deportation game is a very risky game to play, and it's very difficult logistically to get it done. Americans do not like seeing violence against people in the country. [\ldots{}] Hopefully Trump can do this in a savvy way.\footnote{ibid., t. 57:04}
\end{quote}

I think that is about as clear as he can possibly be on this point. He really does endorse this policy, there is no cheeky workaround or re-interpretation possible here. He wasn't being hyperbolic when he called consistent libertarians as bad as communists; in fact what he really wants to say about us is \emph{worse} than that---\emph{by his own admission}. He is only held back by YouTube terms of service. He wants mass deportations done ``by any means necessary,'' and his only hope is that Trump is savvy enough to avoid the public backlash---an odd concern for someone who roots their theory in public opinion. Surely that there would be public backlash implies that the American will is not quite as closed borders as he has previously claimed.

Of course, much cope has been spread over this clip, that he simply wants to dismantle the US empire and that's what bringing them home means. But this isn't actually what he said. He didn't \emph{just} say that he wants to bring them home, there was a second part of his policy proposal. He want's to bring them home \emph{and} set them to the task of ``securing our border'' ``by any means necessary.'' These are separate actions that must be analysed separately. Sure, bringing the US military home is good, that military should not be used to bomb random people on the other side of the globe. But it is very much \emph{not} good to use that military to destroy the last bastion of capitalism on planet Earth.

Consider an analogous policy: I want to take the violent gang members off the streets and place them in senior homes to attack old people. If someone approaches me and says how horrible this would be, what argument could they possibly use? I'm getting them off the streets, aren't I!? We need to live in reality, and in reality this gang violence problem needs to be sorted out. We needn't untangle the package any further than that.

The truth is that this clip is a rare glimpse into what an honest advocacy of immigration restrictions consists of. These restrictions \emph{necessarily} lead to tyranny and \emph{necessarily} require the heavy hand of the state to enforce. You can't eat your cake and have it too; no less can you have restricted immigration without the means \emph{required} to restrict that immigration.

US immigration law is simply different in kind to any other law on the books. Every other law in the Anglo-Saxon tradition of jurisprudence says what you are not allowed to do, leaving everything else as kosher. Immigration law, on the other hand, says what you \emph{are} allowed to do. That is; all immigration is assumed to be prohibited, unless you fit some Soviet-style category where the government grants you permission that absolves you of this original sin. This is why immigration enforcement tends to be way more totalitarian.

This is a system of guilty until proven innocent; and it is no wonder why the Trump administration threw away due process exactly when it came to this area. For every other crime it would be incumbent on the state to prove that there was some harm done; even if it is a nebulous ``crime against society,'' or whatever. With immigration law, you need to prove that you are innocent. That you have a special government permission slip. And you then have to carry your papers everywhere you go or else federal agents might throw you in the back of a van, make up a warrant in the field, and ship you off to a prison camp in El Salvador.

It is because of this \emph{massive} increase in state power--that I repeat \emph{Dave thinks doesn't go far enough}--that you are able to have ICE agents act as the enforcement arm of the state of Israel within US territory. Federal agents detained green card holder Mahmoud Khalil, without a warrant, for expressing views contrary to US foreign policy surrounding Israel. Rümeysa Öztürk, a Turkish student, was snatched from her home by Homeland Security for writing op-eds critical of Israel. ICE says that part of their job is to keep out ideas that cross the border illegally.\footnote{Dan Gooding, \emph{ICE Deletes Post About Stopping 'Illegal Ideas' From Crossing Border}, \url{https://www.newsweek.com/ice-illegal-ideas-border-security-social-media-post-2058217}}

If the state should be restricting foreign-to-domestic immigration per the mystic will of the ``American people,'' should it also stop internal immigration? Could it not be the will of Texans to not allow Californians to immigrate there? We wouldn't want any forced integration, so let's restrict internal immigration as well. But, we needn't stop here. What about the blue cities as against the red country? There could still be a possibility for forced integration, so let's put more restrictions there. But also what if I am a red-blooded Republican and don't want to be forcefully integrated with my blue neighbours; we need to have forced disintegration here too. Heck, let's just have it such that you need government permission to interact with anybody ever, and we find ourselves back at the COVID lockdowns.

When you grant the initial premise that forced integration is to be properly combated with forced disintegration, it never ends. As with all forms of protectionism, you can easily project the reductio case right down to the individual level. Immigration restrictions are, and must be, totalitarian. There's no other way to slice it.
\subsection{Broad, then Narrow}
\label{sec:orgadbde35}
Recognising that there are these legitimate concerns with the closed borders position; the official Mises Caucus policy was at one point a ``big tent'' approach. That is, the Mises Caucus would take up no explicit stance on the borders question, that contentious issues like abortion or immigration should be set to the side to focus on monetary policy and ending the wars. The last surviving evidence of this policy can be found in the fact that there is zero mention of either issue in the entire Mises Caucus Platform.\footnote{\url{https://lpmisescaucus.com/platform/}}

This would certainly line up nicely with how Ron Paul approached the issue. In his many decades-long career in politics, one finds incredibly few instances where he directly talks about borders. And in those few instances he certainly approaches nothing like the current neo-prag stance on the matter:
\begin{quote}
Open borders, it's a more difficult question because under today's circumstances, people see open borders as a very sincere threat to the economy, and to the welfare system, because some aliens can come in and be a burden. In a free society open borders are very very good, just think about how the so-called open borders--which have essentially been open in Canada--people come back and forth all the time and work on both sides and go back and forth, it's not a major problem.

People are more concerned because of the government in Mexico and some of the aliens coming up and I think there's a little bit of racism involved here as well, but in a free society you're always short labour. So people who come up actually give us a service at the same---this is the best foreign aid you could think of. People come up and they work, they take the money back, and they help those families, instead of what we're doing: we tax our people then we give the money to the governments like Mexico and then they rip off the system and they take it and the people never get the money.

So it's the opposite of a socialist approach to helping other people---actually aliens are a net benefit. Under today's circumstances, those who are concerned about aliens coming over, I can understand the concern, but that would not be the case in a libertarian society.\footnote{Ron Paul on Open Borders (1988), \url{https://www.youtube.com/watch?v=e3m41mkA1dY}}
\end{quote}

Notice the way that Dr. Paul handled that question: he took the opportunity to paint the picture of a free society. He leveraged his knowledge of economics to explain to the audience that under capitalism any problems they might be imagining aren't present any longer. That immigration is a net-benefit. He also used it as a jumping-off point to call out the evils of the foreign aid system, and how by letting business actually flourish and hire whomever they want you can ``help'' these people in a purely selfish, capitalistic way, and that this is actually a far larger benefit than any centrally-planned socialist way of dealing with the problem of poverty.

Now, can any person on planet Earth imagine for one second that if Dave Smith is given that same question that he doesn't immediately use it as a chance to countersignal radical, open borders libertarians? I get that I can't expect every commentator to be as adept at navigating these things as Ron Paul, I doubt Dave would say that he is better in that respect; but Dave is taking it in a whole different \emph{direction} to Dr Paul. That is, the problem with Dave's hypothetical response to such a question is not just that he doesn't quite as effectively paint the picture of the free society; it's that he actively takes the border question as an opportunity to shout to the world that radical libertarians believe crazy things that nobody should take seriously. This is an active hindrance to the movement.

As one should expect, the high-time preference with respect to political strategy that has gripped the Mises Caucus leadership made it such that they couldn't follow the policy of Ron Paul, one that they had previously endorsed. They just couldn't keep their big mouths shut and they had to start prancing around calling the open borders libertarians worse than communists, or unconcerned with reality, and so on. The culture war issue of the day had taken hold and proper strategy had to fall to the wayside---they wanted to exacerbate the disagreements on immigration in any way possible.

This emphasis that the mises caucus decided to place on borders was a massive misstep. They should have seen the issue and emphasised that the problem is the welfare state and the civil rights act. They should have used it as an opportunity to denounce those things, rather than endorse immigration restrictions.

As Hoppe cogently points out, immigration becomes problematic with restrictionist trade policies, thus the focus should be on implementing free-trade policies. You know, like calling out Trump's tariffs:
\begin{quote}
Insofar as a high-wage area such as the U.S. engaged in unrestricted free trade, internationally as well as domestically, the immigration pressure from low-wage countries would be kept low or reduced, and hence, the question as to what to do about immigration would be less urgent. On the other hand, insofar as the U.S. engaged in protectionist policies against the products of low-wage area and in welfare policies at home, immigration pressure would be kept high or even raised, and the immigration question would assume great importance in public debate.\footnote{Hans-Hermann Hoppe, ``The Case for Free Trade and Restricted immigration,'' \emph{Journal of Libertarian Studies} 13:2 (Summer 1998): 221-233}
\end{quote}

This terrible decision by the Mises Caucus, to bring the border question to the forefront for temporary culture war wins would, according to Harrison Wells, lead to the eventual split between Mises Arizona and the rest of the Mises Caucus:
\begin{quote}
Affiliates who supported open borders would get nastygrams about how ``we don't talk about immigration'' but affiliates who supported closed borders could post about how good deportations are to their hearts content.\footnote{Harrison Wells via Twitter Direct Messaging.}
\end{quote}

\begin{quote}
They would also hide behind ``the MC has never supported open borders'' whenever the general concept of being principled was brought up.

This was especially true when the person [bringing] it up was me but they did this to other people too.\footnote{ibid.}
\end{quote}

\begin{quote}
I would say ``it is important to be principled and maintain our position as the propertarian wing of the party'' and the counter would be ``we aren't an open borders party'' even if I wasn't talking about immigration but was simply speaking generally.\footnote{ibid.}
\end{quote}

\begin{quote}
This was a common retort to me during the time when I was publicly arguing with the national comms director for saying we need to abandon the NAP.\footnote{ibid.}
\end{quote}

The strategy is laid bare: widen the tent to include both sides of the issue through a public-facing ``neutrality'' stance; only enforce this neutrality for the side you don't want to endorse; now shrink the tent back down to only the side you really wanted all along. They broaden the message, then narrow it back down again. To borrow a favourite term of theirs, this is effectively ``anarcho-tyranny.'' They are perfectly fine with enforcing language norms on the radicals and infighting all day long to keep the message broad-then-narrow; but the instant that a radical attempts to gatekeep the movement and keep out active subverters, they are too idealistic, not strategically minded enough, or too concerned with theory. We need to be accepting of anybody who is vaguely right of center and has a podcast.
\subsection{Neo-Prags as Reactionary}
\label{sec:org610bba3}
This tendency for the neo-prags and the broader conservative movement to shift towards this pragmatism that necessitates tyrannical policies has not come from nowhere. As James Ellias points out,\footnote{James Ellias' untitled lecture on cultural evolution.} the honest, values-oriented individual has been effectively forced to retreat from philosophy by virtue of Immanuel Kant:
\begin{quote}
Good people retreated from philosophy and decided to just live their lives. What this did is it left them more susceptible to faith. I think the better people are more susceptible to faith. Now you might say: ``well, wait a second, faith is awful, you're shutting off your reason! And what do you mean, `the better people!?''' By better people I mean better in the sense of the standard that I set up at the beginning of this lecture: people who are more interested in looking at the facts. Just in everyday life, just as a matter of course, they make more better decisions; they're interested in looking at the facts. Why are they interested in looking at the facts? Because they want to pursue the good, they're not interested in making excuses, they want to get stuff done.\footnote{James Ellias' untitled lecture on cultural evolution, t. 43:00}
\end{quote}

It is clear to see how this tendency leads to pragmatism among these good-oriented people. They have their values and they desperately want to achieve them. They understand on some primitive level that that is what really matters. But they lack any solid philosophical backing, so they are unable to counter arguments from the leftists that there is no such thing as the good, or knowledge, or reality. Thus, an easy alternative is to just do the exact same thing the original pragmatists did in this situation: to invert what it means for something to be true. To take the good as the primary from which truth derives. To say that the true is a species of the good.

This is a completely understandable error for these pragmatist-conservatives to make. Over the past century or so, the general trend has been that the destroyers--that is, those wicked men who killed philosophy and built from its corpse a mangled demon--have been the ones in the drivers seat. This is how Ellias characterises the leftist vs conservative culture war: that the conservatives are mostly good but philosophically confused; which leaves the field wide open for the leftist destroyers who actually have a philosophy to constantly win the argument and destroy every value the conservatives wish to preserve.

After a point they simply have had enough of buying their gay beer, and having their children taught things they consider obviously at odds with basic biological science. They are done with arguments to try and justify their positions because they see the fruits of that path. They see that their values are constantly being eroded and they are powerless to stop it through philosophical means, because they have already abandoned philosophy.

In their view, the only option left to them is to turn to a dogmatic faith in their conception of the good. Those just are The Values™ and they have no ability or intention to justify those values. To defend these values, they of course can therefore no longer argue in favour of them; instead they need to turn to a strongman like Trump to come to their defense.

This was fundamentally a move based upon fear---fear of the left, who have been successfully tearing apart everything they love:
\begin{quote}
They're turning to faith and to a strongman in order to defend their values---the strongman, of course, being Trump, and people like Trump who are going to come after him. In Europe you see people like this too, my understanding is that Viktor Orbán is an example of this, you're going to get all these little orange men popping up all over, [\ldots{}] who are willing to just basically break the rules in order to give these people what they want. And they want it for sort of legitimate reasons: they can see their values being destroyed bit by bit by the left and they don't know how to stand up for them in a way that will actually preserve what the enlightenment has achieved. And so why not turn to this strongman?\footnote{James Elias' unnamed lecture on cultural evolution, t. 57:10}
\end{quote}

To borrow a term that the leftists love: this is fundamentally a ``reactionary'' attitude that they are taking. The leftists will attack some dogmatically held values--some sacred cow--that the conservatives wish to preserve. Through just a basic almost-instinctual fear response; the conservatives will put all of their focus on responding to this new attack.

They will see all of these terrible things that the leftists are doing--destroying everything the conservative wishes to pursue--and the character of their response is ``I want anything but \emph{that}! I'll take whatever else except for leftism.'' The problem with such fear-based responses in turning to the strongman to stamp out the immediate problem is that they simply cannot serve to actually preserve the values that must be preserved. The reactionary attitude leads only to further philosophical degredation.

Because when you say ``anything but that'' you really are willing to accept anything but that. This gives an open license to any tyrant who wishes to come along and provide an alternative. The only way to preserve objective values is to adhere to and promulgate a philosophy that has a rational, this-worldly justification for those values. If you try to come up with some other-worldly justification you end up simply creating the very higher-order rationalisations that made the leftist destruction of philosophy possible in the first place.

That is, when you act out of fear in this manner, your actions become random with respect to values. You may be able to avoid that one evil that you have within your focus, but in the process you open the door for a litany of additional evils.

This is why the fearful response to the border crisis from Dave wherein he advocated an even-more extreme policy than Trump inevitably leads to unforeseen consequences---that is, unforeseen by the neo-prag. We principled libertarians were able to easily foresee how militarising the border would go part-in-parcel with a greater war-hawk attitude from Trump; but Dave and his compatriots were unable to consider this obvious truth. They were blinded by their high time preference pragmatic attitude and unable to think through the broader philosophical implications. They have no philosophy, so they are presented only with the options of engaging in this reactionary dogmatism, or falling to the dark side and becoming a leftist.

This is not challenging to pick out in the Dave Smith vs Andrew Wilson debate:
\begin{quote}
Andrew: Can you justify why two consenting adults who are brother and sister having sex is morally wrong? [\ldots{}]

Dave: For the same reason you would think: I think it's against nature, it's evil, it just seems horrible. Do you have a better answer for why it is wrong?

Andrew: Yes, divine command.

Dave: OK, then I also say divine command.\footnote{The Crucible, \emph{Andrew Wilson vs Dave Smith: Is Libertarianism better than Christian Populism?}, \url{https://www.youtube.com/watch?v=YqIaiQ-aK\_s}, t. 31:44}

{[}\ldots{}]

Dave: I would say: I grew up an atheist, [\ldots{}] and then I had some experiences as a young adult that made me think about god and kind of get in touch with him, and it was really when I had my first child when I really found god. So I would agree with you that there is something from god that instructs me as to what is right and wrong in this world. However, before I ever believed in god, yeah there was just a deep feeling inside of me that that would have been very wrong.\footnote{The Crucible, \emph{Andrew Wilson vs Dave Smith: Is Libertarianism better than Christian Populism?}, \url{https://www.youtube.com/watch?v=YqIaiQ-aK\_s}, t. 37:18}

{[}\ldots{}]

Super Chat: Dave, [by] what objective standard do you ground moral laws such as the NAP or anything for that matter? I know you believe in God, but if that's just general Theism, by what standard do you ground any objective morality?

Dave: I don't think me or Andrew had a better answer for that; I've always just thought--before I was a believer in god--that it was kind of a silly argument that theists made that you can't have morality without god. Since I've found god, I've agreed that yeah, you probably can't. There really is no other way to get to morality without belief in a creator. [\ldots{}] Since I had my first kid I've just kind of known, it's just been more obvious that most of the things that are just kind of deep-rooted in us as inherently wrong or right come from that.\footnote{The Crucible, \emph{Andrew Wilson vs Dave Smith: Is Libertarianism better than Christian Populism?}, \url{https://www.youtube.com/watch?v=YqIaiQ-aK\_s}, t. 1:38:08}
\end{quote}

That is about as explicit as we could hope him to be on this front. He is against incest because it just ``feels wrong,'' and is forced to adopt the first philosophical argument for his feelings that was provided by Andrew. This argument, of course, just being purely taken on faith by Dave. It is after he had his first child that the values he already implicitly held to became utterly apparent to him and achieved a brand-new urgency. At this point he is pushed into the reactionary state of total fear. He has no objective standard to ground his moral laws, as he admits, and so he tells us straight up that he just always had those values and turned to dogmatism as a basis to protect them. It then should be no surprise to you, dear viewer, that Dave has since abandoned any basis required for libertarianism and endorsed some of the most authoritarian legislation one could imagine.

This fear-based morality that the Mises Caucus has adopted will necessarily spell their defeat. The philosophically inclined will take advantage of them at every turn, and they will never develop the weapons to effectively fight for their values.
\section{On Yellow-Washing, and Greening Out}
\label{sec:orgf6b7e17}
The Mises Caucus had their first big victory at the 2022 Libertarian National Convention, in what was called the ``Reno Reset''---an attempted ousting of any pragmatism in the party in favour of a new radical direction. Two years later, at the 2024 convention, would come one of their greatest defeats. They went from achieving landslide victories for every single one of their candidates, to barely scraping by and not even getting their presidential nominee on the ticket. What happened to these political mastermind Machiavellian geniuses, and why did they fall so far so fast?

Of course, Dave Smith was pegged to be their champion, their chosen candidate. He had all but explicitly agreed before pulling out at the last minute to focus on his family. This had the effect of basically removing any chance for the Mises Caucus to put forward a decent candidate. As former chair Angela McArdle Explains:
\begin{quote}
We had 24\%-26\% in 2018. We had \textasciitilde{}40\% in 2020. We would have come up with about 60\% in 2022, and we came in with 70\% instead. What happened when Dave stepped in and said he was going to run for president is the mission completely wrapped itself around that. And instead of just decentralised revolution [\ldots{}] it became the Dave campaign. So when everything wrapped around Dave that meant there was a lot at stake with whatever Dave Smith did, and for several reasons [\ldots{}] Dave decided not to run, and when that happened it created a huge vacuum, a huge shockwave of disappointment, and general chaos. And it was really challenging to reclaim that.\footnote{Michael Malice, \emph{``YOUR WELCOME'' with Michael Malice \#314: Angela McArdle}, \url{https://www.youtube.com/watch?v=EYIk9KbnqLE}, t. 29:04}
\end{quote}

\begin{quote}
Dave Smith was our chosen candidate--and he's owned this publicly, so I'm not dragging Dave any more than he has dragged himself--he was gonna run, he announced it publicly in a very teasing way. But multiple times we had conversations about it behind the scenes, and he decided he could not after a year of us being deadlocked and one of the things that happened as a result of Dave announcing is other people who might have shown interest, Justin Amash for one, decided not to run. Jo Jorgensen another. People were like: ``Dave Smith is running, we're not gonna be able to compete, we're not gonna go there,'' and then they just found other things to do and they had stuff going on in their lives. So when the Dave Smith rug got pulled out from under us it left a giant vacuum.\footnote{Angela McArdle, \emph{The Much Anticipated Endorsement Announcement}, \url{https://www.youtube.com/watch?v=bGkWugKoKWU}, t. 2:46}
\end{quote}

Now, I want to be clear here, I don't give a single shit that Dave screwed over all these pragmatic ``Realpolitik'' believers who had already at this point abandoned any pretence of radicalism. Certainly, spending time with his family is vastly more important than running a presidential campaign; and as Dave himself tells us, the Libertarian candidate for president should be a libertarian---so on that alone I'm glad he didn't run.

Luckily for the Mises caucus, the clouds would part and present them the perfect candidate. A rough-and-tumble former-professional-communist, about whom the Mises Caucus founder had such kind words as ``he was the only choice,'' Michael Rectenwald had arrived and was ready to shake things up.

And, I want to be clear about the timeline here. By ``former-professional-communist'' I don't mean that when he was 18 he joined the Black Panthers to get laid or some shit; I mean that he was a self-described Marxist and ``left-communist''\footnote{Diamond Naga Siu, \emph{Q\&A with a Deplorable NYU Professor} (Washington Square News), \url{https://nyunews.com/2016/10/24/qa-with-a-deplorable-nyu-professor/}} well into his long career as a Liberal Studies professor at NYU. That is, he was engaged in his undergraduate studies in the late 70s and early 80s, and was in love with the works of communist poet, Allen Ginsberg.\footnote{Amber Frost, \emph{A Dangerous Minds exclusive: Previously unpublished interview with Allen Ginsberg}, \url{https://dangerousminds.net/comments/previously\_unpublished\_interview\_with\_allen\_ginsberg/}} Michael was a true believer and would dedicate the following four decades of his life to bringing about the revolution that Marx had promised to him.

It was only after the left tried to cancel him that he started to re-think his positions. As he puts it:
\begin{quote}
The left did more for my conversion than anybody. I think it was a trauma, I suffered a trauma, and that trauma actually opened up my eyes. They did for me what I couldn't do for myself\footnote{BlazeTV, \emph{Former Marxist Turned Trump Voter On How to Defeat Globalists | Zero Hour | Ep 9}, \url{https://www.youtube.com/watch?v=uKiFXHOpuxk}, t. 12:15}

{[}\ldots{}]

I was a Marxist in September and I voted for Trump in November.\footnote{BlazeTV, \emph{Former Marxist Turned Trump Voter On How to Defeat Globalists | Zero Hour | Ep 9}, \url{https://www.youtube.com/watch?v=uKiFXHOpuxk}, t. 14:43}
\end{quote}

I think that we as a rational and highly radical movement have a certain duty when it comes to being careful about who we select as thought leaders. Someone who was a career Marxist right up until the point when the left tries to cancel them leading them to turn coat, is not someone who should be immediately embraced with open arms. In this state, Rectenwald is engaged in a form of motivated reasoning; he ``left the left'' because of how meanly he was treated, not because he had fully developed his brain and saw the evils that his comrades were advocating. He didn't read some super-convincing argument to get converted; he got called names on the internet. We should be suspicious of anybody who is a supposedly-former communist attempting to enter our movement; and we should be doubly suspicious of those who were career-communists that only left the fold when the machine they were building turned against them.

The correct attitude to people like Rectenwald is almost total suspicion. He should have had to go through the ringer and prove to everybody else that he truly had changed, that he utterly renounced his former affiliations. as Peikoff states the point:
\begin{quote}
The general principle here is: truth implies as its cause a virtuous mental process; falsehood, beyond a certain point, implies a process of vice. [\ldots{}]

It is possible for a man to embrace an idea blindly, on faith from others or simply by his own whim, without the effort of understanding or integrating it. In such a case, the idea, no matter what its content, reflects negatively on the individual. For Objectivism, an idea thus embraced is not “true” (or “false”). In relation to such a mind, the idea is without cognitive status; it is the arbitrary, and is analogous to the sounds emitted by a parrot. The true qua true, by contrast, does imply a process of understanding and integration, and therefore some degree of effort, focus, work. [\ldots{}]

Now we must note that falsehood does not necessarily imply vice; honest errors of knowledge are possible. But such errors are not nearly so common as some people wish to think, especially in the field of philosophy. In our century, there have been countless mass movements dedicated to \emph{inherently dishonest} ideas --- e.g., Nazism, Communism, non-objective art, non-Aristotelian logic, egalitarianism, nihilism, the pragmatist cult of compromise, the Shirley MacLaine types, who “channel” with ghosts and recount their previous lives; etc. In all such cases, the ideas are not merely false; in one form or another, they represent an explicit rebellion against reason and reality (and, therefore, against man and values). If the conscientious attempt to perceive reality by the use of one’s mind is the essence of honesty, no such rebellion can qualify as “honest.”

The originators, leaders and intellectual spokesmen of all such movements are necessarily evaders on a major scale; they are not merely mistaken, but are crusading irrationalists. The mass base of such movements are not evaders of the same kind; but most of the followers are dishonest in their own passive way. They are unthinking, intellectually irresponsible ballast, unconcerned with logic or truth. They go along with corrupt trend-setters because their neighbors demand it, and/or because a given notion satisfies some out-of-context desire they happen to feel. People of this kind are not the helplessly ignorant, but the willfully self-deluded.

Even in regard to inherently dishonest movements, let me now add, a marginal third category of adherent is possible: the relatively small number who struggle conscientiously, but simply cannot grasp the issues and the monumental corruption involved. These are the handful who become Communists, “channelers,” etc. through a truly honest error of knowledge. Leaving aside the retarded and the illiterate, who are effectively helpless in such matters, this third group consists almost exclusively of the very young --- and precisely for this reason, these youngsters get out of such movements fast, on their own, without needing lectures from others; they get out as they reach maturity. Being conscientious and mentally active, they see first-hand what is going on in their movement and they identify what it means; so their initial enthusiasm turns to dismay and then to horror. (Andrei in \emph{We the Living} may be taken as a fictional symbol here.) The very honesty of such individuals limits their stay in the movement; they cannot tolerate for long the massiveness of the evil with which they have become involved. Nor, when such youngsters drop out, do they say to the world belligerently: “Don’t dare to judge me for my past, because my error was honest.” On the contrary --- and here I speak from my own personal experience of honest errors that I committed as a teenager --- the best among these young people are contrite; they recognize the aid and comfort, inadvertent though it be, which they have been giving to error and evil, and they seek to make amends for it. They expect those who know of their past creeds and allegiances to regard them with suspicion; they know that it is their own responsibility to demonstrate \emph{objectively} and across time that they have changed, that they will not repeat their error tomorrow in another variant, that their error was innocent.\footnote{Leonard Peikoff, \emph{Fact and Value}, \url{https://peikoff.com/essays\_and\_articles/fact-and-value/}}
\end{quote}

It is on this front that Michael Rectenwald should have been nowhere near any prominent or leadership position in the libertarian movement. The man was on an active war against reality and the good for basically his entire adult life; and the second said war was turned against him \emph{that} is when he finally decided to change sides.

Given these circumstances, and the breakneck speed involved in getting bullied by the left in February,\footnote{\url{https://web.archive.org/web/20170302003856/https://twitter.com/antipcnyuprof/status/834426977616527360}} followed by supposedly being a full-on Misesian anarchist by fall,\footnote{Michael Rectenwald, \emph{How a Marxist of Twenty-Five Years Became a Misesian Libertarian}, \url{https://mises.org/mises-wire/how-marxist-twenty-five-years-became-misesian-libertarian}} it is incumbent upon us to leave this guy to the fringes.

It is the same as the case of leftist feminists who deliberately created the world we live in--trampling on the rights of everyone in their path--who now cry out about how their free speech is being violated for speaking out about the transgender issue. None of these people who have spent a career advocating for the Leopards Eating People's Faces Party should be granted a prominent role within our movement the moment that the leopards eat \emph{their} face.

It really is no wonder that Rectenwald has started returning to his socialist roots now that the dust has settled. He has bounced around the political compass once more to land on Jew-hatred being his new Categorical Imperative. Allow me to go over some key points in the platform for his proposed ``Anti-Zionist America Party.''\footnote{\url{https://x.com/RecTheRegime/status/1945487226357305768}}

He wishes to ``dismantle the military-industrial complex by cutting the defense budget by at least 50\%, redirecting funds to border security, infrastructure, and debt reduction.'' Now, let me be clear about the purpose of a party platform. This is supposed to be a \emph{statement of principles} that he and the fellow members of the party are to endorse. This is not the place to expound on policies that you think would be better than the status quo; but the \emph{ideals} that one is striving towards achieving.

Given this, if the principle is to ``dismantle the military-industrial complex,'' in what sense is cutting the budget by half a realisation of that principle? Is that not the \emph{very definition} of a half-measure? Consider if the slavery abolitionists put forward as their statement of principle, that imports of slaves be cut down to half of the current rate. Would you consider this person a true abolitionist in any sense of the word?

Further, redirecting these stolen assets towards ICE and the Department of Transportation is certainly \emph{not} the libertarian position. If this is his ideal, in no sense can he be described as such. To be clear; the libertarian position on the military industrial complex is that the United States should have no standing military; and the libertarian position on taxes is that the state should not collect taxes---not that those taxes go towards Rectenwald's favourite pet projects.

He continues with his ``Pro-Free Speech and Anti-Censorship'' plank as follows: ``Repeal laws like the PATRIOT Act and FISA amendments that enable government surveillance. Protect online platforms from corporate censorship by enforcing antitrust laws against tech monopolies. Criminalize attempts to suppress speech on topics like foreign policy, history, or government criticism, including Holocaust revisionism or discussions of Zionist influence.''

It barely necessitates enumeration on the ways in which this is an anti-libertarian stance. Certainly, state legislation such as the PATRIOT Act and FISA should be repealed, that is the easy part. But in no sense is it in any way libertarian to have a big government come in to break up companies or enforce anti-discrimination laws on people who do not wish to platform views they disagree with. It is any man's right to exclude Holocaust deniers from his own damn property.

The party position on borders is, to nobodies surprise, also authoritarian garbage: ``Impose a complete halt to all immigration for at least five years to allow for national reassessment and security upgrades. Secure borders to prevent exploitation by foreign interests, while deporting existing illegal entrants and reforming the system to prioritize American sovereignty.''

I think it bears no repetition in this video on how draconian the policy of a complete sealing of the border is; but I do want to draw your attention to the distinctly Marxian language he uses in defense of that policy. He is concerned about ``foreign interests'' ``exploiting'' people. What exactly could that possibly mean from a professional leftist if not that a private company hiring individuals to work somehow constitutes exploitation of the workers?

He continues in advocating that the government centrally plan international trade through the use of tariffs and the promotion of ``energy independence through domestic sources.'' What does it mean for a government to \emph{promote} some industry if not subsidies and/or some other regulation. This man has clearly abandoned any pretence of being a Misesian at this point.

I bring up just how far Rectenwald has strayed back into his old socialist habits to highlight a major problem with people like McArdle. Much has been said regarding her strategy of leveraging the spoiler effect to obtain concessions from Trump such as the freeing of Ross Ulbricht and getting a libertarian on the cabinet. Of course, I have absolutely no problem with using the spoiler effect in this way, that is not an inherently bad idea. The problem with Angela and others at the Mises Caucus, is that this strategy is packaged together with a process of ``yellow-washing'' our ideological enemies, mostly Republicans and other ``dissident rightists.''

This is obviously apparent in the case of Rectenwald; with McArdle spinning his former occupation into a positive, calling it ``an incredible redemption story arc,'' rather than justly treating him with suspicion. But she extends the yellow-washing agenda even to the arch-socialist that is Donald Trump, calling him a ``Libertarian POTUS.'' This is either Nick Sarwark-esque ``party-over-principles'' pragmatism, or a complete bait-and-switch. We were promised libertarians; irrespective of party membership:

\begin{quote}
Donald Trump says he's going to put a libertarian in a cabinet position. He came out and spoke to us. He said he's a libertarian. He has basically endorsed us, and so in return, I endorse Chase Oliver as the best way to beat Joe Biden. Get in loser, we are stopping Biden. That's what I think. That's what I think this campaign is about.\footnote{Angela McArdle, \emph{The Much Anticipated Endorsement Announcement}, \url{https://www.youtube.com/watch?v=bGkWugKoKWU}, t. 24:13}
\end{quote}

In no sense is Donald Trump anything approaching a libertarian; but perhaps he at least put a libertarian in a cabinet position. I certainly can't see anybody; Trump instead chose a number of MAGA sycophants, and Angela is seemingly honest enough to agree that this concession did not come to fruition:

\begin{quote}
Interviewer: If it's just some minor official in some minor government post, or if it's RFK Jr. who is right now very proximate to Trump and who I agree with on some things [\ldots{}]. But also has a long history of supporting environmental regulations and health regulations in ways that are not libertarian; so paradoxically if the libertarian choice is for him in a role which he does not align with libertarianism, we can't call that a win, right?

Angela: No.\footnote{The Hill, \emph{Trump APPOINTING Libertarians To Administration?! Interview with LP Chair Angela McArdle}, \url{https://www.youtube.com/watch?v=R4r1OaXBAyM}, t. 6:07}
\end{quote}

So by Angela's statement here; we \emph{cannot} call RFK a win or a libertarian pick because he has been placed in one of the two named areas, that is healthcare, where he is decidedly not libertarian. This makes her statements from later in that very same month when addressing Dave Smith's twitter audience, quite puzzling:
\begin{quote}
I think the goalposts have shifted considerably in the cabinet conversation. I want to remind everyone that Trump never promised an entirely libertarian administration, those of you pointing and screeching about every single pick. Like it's not correct. I have repeatedly said we will only get one, only one cabinet position and positions underneath that.

No one promised you a better GOP either. I said we will get one and we will have to keep fighting for more every election cycle. And I also said that if Trump only keeps one promise, if he only frees Ross, it will still be a win and the most progress we have ever made at the federal level.

Now I will say that Robert F. Kennedy Jr. joined the party as a lifetime member. He joined as a lifetime member. Kennedy is our guy, like it or not. He's putting our people into the administration. The man joined the party. He sought our nomination. He didn't get it. And after that, he joined as a lifetime member. After he lost our nomination.

He's gone to bat for me and for all of us with the incoming administration, I'm going to bat for him right now. And he had leverage, by the way, because he dropped out and endorsed Donald Trump. That's why he and his team are at Mar-a-Lago right now, going to bat for libertarians against the bad actors in the Senate and AIPAC and all the other people who are pressuring to try to push us out.

And I am so grateful for him for the work that he's doing and the ground he's forging for us. And this is one of the reasons also that I have so aggressively defended the LNC's decision to engage in a joint fundraising agreement with him, which we are still doing.

I have, by the way, already done reference calls for some of the libertarians who have applied to be in the administration and filled out many applications and referrals myself. That is happening, and a lot of those picks are up to Kennedy himself.

Go ahead and scream. Go ahead and scream that Kennedy isn't the one true libertarian. It's not good enough that he displayed moral courage during a national crisis where our two most recent LP presidential candidates were wearing masks and talking about social distancing, outdoor-only events, taking the Fauci-owchie, you are all possessed.

And I want to say, I'm speaking to my critics right now. You suffer from ideological paralysis. You are possessed by a very unproductive demon, shackled down in paralysis and unable to move because you can't find the purely libertarian path forward.

That's not how politics works.

Who here has ever tried to get a job in the administration? Who here helped me vet candidates for a cabinet position? Who made calls and gave me names? Only a few of you. Thank you to those of you who did.

Do not complain about the cabinet position if you do not help us try to achieve it. At least don't complain about the work that I've done. If you sat on the sidelines and sulked, don't complain. I don't want to hear it.\footnote{Angela McArdle speaking on Dave Smith's twitter space, \emph{Did I sell out my principles by supporting Trump?}, \url{https://x.com/ComicDaveSmith/status/1860865805006590365}, t. 20:43}
\end{quote}

This would make it seem very much the case that she considers Kennedy to be ``our guy''---that is: that Kennedy is the libertarian in the cabinet. But when talking to a more normie audience in her interview with \emph{The Hill} she is perfectly able to accept that Kennedy is a socialist specifically in the area he was selected to run. So in other words, when Angela is talking to non-libertarians she completely drops the yellow-washing, she agrees that in no way can RFK as Health Secretary be counted as a libertarian win; but the moment that she is speaking to a libertarian audience, she spends all of her time expounding on the virtues of RFK and countersignalling any inkling of radicalism.

She also gives us there her standard for what makes RFK ``our guy,'' namely that he has ``moral courage.'' She would redouble this standard later in the space when speaking to lolcow-in-chief, Jeremiah Harding:
\begin{quote}
Jeremiah Harding: And I recognize that you're a political party, but you acted pretty much not like that when people like RFK Jr. can pay a little bit of money and a little bit of lip service while having the record he does in the like statist factions that he's been in and the actual things that he has participated in and still get a seat at the table in terms of being our guy.

No, he's not.

RFK is not my guy and he's not the guy of anybody with any actual principles. But you're saying he is because he gave you some money and because he's giving you some lip service.

Angela McArdle: He doesn't have to be your guy, but there are a lot of people who are super excited about having someone with moral courage to go after Fauci and go after the COVID regime.\footnote{Jeremiah Harding and Angela McArdle speaking on Dave Smith's twitter space, \emph{Did I sell out my principles by supporting Trump?}, \url{https://x.com/ComicDaveSmith/status/1860865805006590365}, t. 1:20:59}
\end{quote}

What exactly does it mean for one to have Moral Courage? Just sticking to their morality---\emph{any} morality? She elaborates in a livestream:
\begin{quote}
I think that we're really fighting over moral courage.\footnote{Angela McArdle, \emph{Moral Courage And The Libertarian Party}, \url{https://www.youtube.com/watch?v=Uvuht8KfLu4}, t. 0:45}

{[}\ldots{}]

I think what people are asking for, essentially, is moral courage. Moral courage is the ability to act on moral reasons, even when there is a risk of negative consequences.\footnote{ibid., t. 1:37}
\end{quote}

So, RFK is ``\emph{our} guy,'' because he has the ability to act on moral reasons, even when there is a risk of negative consequences. Nowhere in her definition does she specify that those moral reasons must be \emph{libertarian} moral reasons. Indeed, Lenin surely had moral courage by this standard; and therefore Lenin must also be our guy! This is not only a terrible conceptual category to use in an ethical evaluation (the proper alternative is between the thinkers and the second-handers, not Angela's standard); but it is a doubly terrible standard to use in judging the strategic advantage we have been afforded by a certain person's victory.

That RFK has moral courage, that he is willing to say what he thinks radically, and boldly, is an absolutely terrible thing for us. Because RFK is a raging socialist. He should be meek and frightened to alienate voters. Such meekness would be far more advantageous to us, because meekness fails where radicalism succeeds. People are not ``our guy'' simply because they are radicals---they must be radicals \emph{for capitalism}. Moreover, what this highlights is that the moral courage standard runs completely counter to pragmatism. Pragmatism tells man to \emph{not} act on moral reasons---that's the whole point! You can't extol the pragmatic virtues of alignment with RFK on the basis that RFK rejects pragmatism and this makes him a good person. If it is good to reject pragmatism, then you should not be making any of your arguments based upon it.

It is no surprise, then, that Angela occupies a great deal of her praxis with libertarian countersignalling, whilst making sure to cozy up to every former-professional leftist and current environmental-communist that comes her way. If her standard is this warped nonsense, what else could possibly come about?

This standard that forces Angela to side with our enemies, and lie to libertarians about those enemies allegiances, can be seen elsewhere within the Mises Caucus. The IP-socialists\footnote{Stephan Kinsella, \emph{Colorado IP Socialists Trying to Amend LPCO Platform to Include IP}, \url{https://c4sif.org/2025/06/colorado-ip-socialists/}} in the Colorado LP, which is pretty much exclusively Mises Caucus-controlled, have a whole yellow-washing campaign with their ``Project Wild Horse.'' The plan was to endorse RFK Jr for president, as stated by the LPCO chairwoman, Hannah Goodman:
\begin{quote}
The Strategy of Colorado [\ldots{}] is to radically reshape the political landscape of Colorado. I'm in a blue super majority. We have literal communists in our state legislature. [\ldots{}] How the fuck do we break up this stranglehold? This play has made Colorado a swing state now. No one is saying RFK is a pure libertarian. Chase by his own admission said he's only broadly libertarian. I'd argue RFK could be considered broadly libertarian then. RFK is polling higher than Biden in Colorado. Colorado doesn't care about internal party politics. We care about taking down the state, we care about breaking up the establishment. [\ldots{}] I don't want to grow members to the party till I grow voters. My [Key Performance Indicator] is gaining more votes for candidates.\footnote{Mises Caucus Discord Server}
\end{quote}

I think that she has slipped the mask here. She says that ``we care about taking down the state, we care about breaking up the establishment.'' To these sort of people, taking down the state involves being on the anti-establishment side of the culture war. Again, they are fundamentally reactionary: they see only the current establishment, put all of their focus on that, and will take whatever plan they come across to put a wrench in the works; even if the plan ends up creating a whole new establishment. That is, anti-establishment is conflated with anti-state, blinding these people and making them endorse active socialists, so long as those socialists are on the right side of the culture war.

It is quite telling that her ``key performance indicator'' is getting more votes for candidates, regardless of who that candidate is. Heck, just endorse whichever two-party candidate is most likely to win in every election and you have successfully achieved a far higher performance than ever before. This sort of votes-over-principles, Nick Sarwarkian pragmatism was once a joke for the Mises Caucus to laugh at. It is now official Mises Caucus policy.

Again, using the spoiler effect and trying to turn Colorado into a swing state is fine, that is not the issue; the issue is that you need to have a strong radical basis to even evaluate when to spoil or not in the first place. These pragmatists simply don't have that. They are cursed to treat every future election as the most important election of our lifetimes forever into the future because the election is the issue of the day. Their pragmatism demands that they put aside any principles and just deal with the election as it comes. This leads to way less spoiling than there should be.

To borrow Angela's term; I do not think that high-time preference pragmatists have anything approaching the moral courage to ever actually follow through in spoiling when it is required. Would any one of these people put aside their culture war grievances to actually punish some major party candidate who failed to follow through on their concessions? I see absolutely no evidence that such is the case.

This is the real problem with yellow-washing, it blurs what should be hardcore moral lines and makes it impossible to distinguish friend from foe; and that poisons any analysis from there on. Fundamentally, it is simply the error of contradiction: it is treating not-A as A; it calls for identifying non-libertarians as libertarians. ``Why shouldn't the Libertarian Party endorse RFK? After all, he's a libertarian, not a raging socialist; the tent has expanded and we need to ignore these autistic radicals!'' ``Who cares that this politician didn't follow through on those concessions he promised last time? He has since had a redemption arc, and is no longer the career-socialist he once was!'' ``Why not advocate for martial law in the United States? This was the principled libertarian stance all along! You people are just too concerned with theory.''

The Machiavellian manoeuvring of the political masterminds behind the Mises Caucus would come to a head at the 2024 LNC. As discussed, they achieved vastly diminished results as compared to the Reno Reset. As Dave explains, he went through the first three days of the conference confident that Rectenwald would be the nominee,\footnote{Dave Smith, \emph{Libertarian National Convention 2024}, \url{https://rumble.com/v4yhk7e-may-30-2024.html}, t. 35:40} so what happened?

The first nail in Rectenwald's coffin, is alleged by McArdle to be the seating of delegates who were not supposed to be able to vote:
\begin{quote}
The haters run a very active campaign of messaging to the neutrals to get them to pull in their direction on a few key points. And one was to seat delegates improperly. Which means: people showed up, they were not delegates, they were not qualified to participate in our proceedings, and it takes a 7/8ths vote to seat them, and I ruled that, and they overturned my ruling.

Here's the trick: it's quite a bit of work, but it only takes [\ldots{}] a majority to overturn my ruling. So there's a cheat code. And if you make it sympathetic, ``oh my gosh, Angela and the Mises Caucus, they don't want these people who travelled this far to participate, that's so unfair.'' And you just get 51\% of the room to agree with you, then you can seat people who hate us. So our numbers drop.\footnote{Michael Malice, \emph{``YOUR WELCOME'' with Michael Malice \#314: Angela McArdle}, \url{https://www.youtube.com/watch?v=EYIk9KbnqLE}, t. 37:57}
\end{quote}

Now, indeed, I am not interested in digging through Libertarian Party bylaws, so I will take Angela's word for it that these people were improperly seated. If her allegations are correct, that is indeed a very scummy, dishonest thing to do. But, why am I hearing this from these mastermind Realpolitik Machiavellians? Isn't this just indicative of her haters' will to power or whatever? If we aren't being radical libertarians anymore, in what possible way might we complain?

Regardless, the convention pushed forth through many arduous rounds of voting, eventually narrowing down to three candidates: Chase Oliver, Michael Rectenwald, and Mike Ter Maat. As a bonus for these three top picks, they were each given the opportunity of a lifetime. Trump was to speak at the convention, and they would be allowed a press conference immediately after to provide any rejoinders they had.

This is where the second nail would come in, because right before what is probably the single most important media appearance of his entire life, Rectenwald decided it would be a great idea to consume a THC edible, in what would become known as ``gummygate.''

The audience became immediately aware of this fact when he tried to stumble through his answers; and his fellow candidates picked up on this too. Particularly, Ter Maat started publicly flirting with a Chase Oliver campaign right there on stage after Rectenwald apparently got too high to continue, and had to leave:
\begin{quote}
I thought that was a good answer!\footnote{News2Share, \emph{Libertarian presidential candidate press conference following Trump speech}, \url{https://www.youtube.com/watch?v=26zILo33hXI}, t. 15:30}

{[}\ldots{}]

Chase and I look like brothers, damn it!\footnote{News2Share, \emph{Libertarian presidential candidate press conference following Trump speech}, \url{https://www.youtube.com/watch?v=26zILo33hXI}, t. 16:44}

{[}\ldots{}]

Some day, eventually, Chase will make a mistake.\footnote{News2Share, \emph{Libertarian presidential candidate press conference following Trump speech}, \url{https://www.youtube.com/watch?v=26zILo33hXI}, t. 17:10}

{[}\ldots{}]

I believe it will be me as the nominee, but no matter who it is, if it's Chase, if it's anyone else [\ldots{}]\footnote{News2Share, \emph{Libertarian presidential candidate press conference following Trump speech}, \url{https://www.youtube.com/watch?v=26zILo33hXI}, t. 18:18}

{[}\ldots{}]

{[}Ter Maat is seen vigorously clapping for Chase's answer]\footnote{News2Share, \emph{Libertarian presidential candidate press conference following Trump speech}, \url{https://www.youtube.com/watch?v=26zILo33hXI}, t. 20:20}

{[}\ldots{}]

{[}Ter Maat gestures over to Chase shouting out a crowd-riling answer during his speaking segment]\footnote{News2Share, \emph{Libertarian presidential candidate press conference following Trump speech}, \url{https://www.youtube.com/watch?v=26zILo33hXI}, t. 20:59}

{[}\ldots{}]

{[}Ter Maat gestures to Oliver] No offense, because I think you would be a great running mate for Donald Trump.\footnote{News2Share, \emph{Libertarian presidential candidate press conference following Trump speech}, \url{https://www.youtube.com/watch?v=26zILo33hXI}, t. 21:37}

{[}\ldots{}]

Oh I thought Chase nailed it.\footnote{News2Share, \emph{Libertarian presidential candidate press conference following Trump speech}, \url{https://www.youtube.com/watch?v=26zILo33hXI}, t. 22:50}

{[}\ldots{}]

{[}etc.]
\end{quote}

Prior to this point, Ter Maat had given assurances to the Mises Caucus that he would not endorse Chase Oliver; but in the next round of voting, when Ter Maat was eliminated, he did the unthinkable, and joined up with Chase as his Vice Presidential candidate---forming the final nail in Rectenwald's coffin. And, of course, the screeching from the Mises Caucus would reach a fever pitch at this point:
\begin{quote}
I went into Saturday night feeling pretty good. Feeling pretty good about where we were and thinking Michael Rectenwald was the favourite to be our presidential nominee. [\ldots{}] He had a bad moment on stage before Donald Trump.\footnote{Dave Smith, \emph{Libertarian National Convention 2024}, \url{https://rumble.com/v4yhk7e-may-30-2024.html}, t. 35:34}

{[}\ldots{}]

And so it comes in and Rectenwald is in first place on the first ballot. [\ldots{}] It went I think 8 or 9 rounds, and Rectenwald was in first place the entire time. So like, Michael Heise, who is a god-damn genius, had like mapped all of this out [\ldots{}] and he was completely right about the whole thing. And he goes: ``the last thing that is going to happen is this guy Mike Ter Maat is gonna come off the board,'' and he's already assured Michael Heise and Rectenwald multiple times, he's promised them, that he knows Chase winning is the worst thing for the party. And so he will endorse Rectenwald when he comes off, and we're like: with that endorsement, we're pretty much a lock to get the nomination. [\ldots{}] It finally comes down to Chase, Rectenwald and Ter Maat comes off the board, and he goes over and tells Chase that if Chase makes him his VP, he'll support Chase. So he lied to Heise and to Rectenwald. I mean, listen, it's politics, this shit happens, he did what was ultimately in his interests, to the detriment of the party, and to the detriment of keeping your word. [\ldots{}] It's not something I would ever do. I think if you have an ounce of integrity you don't lie like that, but whatever.\footnote{Dave Smith, \emph{Libertarian National Convention 2024}, \url{https://rumble.com/v4yhk7e-may-30-2024.html}, t. 37:10}

{[}\ldots{}]

It finished with Rectenwald at about 45\%---he got all of our people, but he didn't get anybody else. And so that's where we are: Chase is the presidential nominee, the Mises Caucus still controls the LNC. But this guy is gonna run for president with Mike Ter Maat, the lying cop, as his vice president.\footnote{Dave Smith, \emph{Libertarian National Convention 2024}, \url{https://rumble.com/v4yhk7e-may-30-2024.html}, t. 39:10}

{[}\ldots{}]

I'm really trying throughout this, one not to be a sore loser, and two to really take the blame for this. Because I do think this is on me. And I see people putting blame on Heise and Rectenwald and I just don't think that's right. I don't think these guys did anything wrong. I was the one who was planning on running and made that known, and I got people excited for that, and I was the one who pulled out. I didn't exactly promise them I was going to run, but I certainly led people to believe that because I was planning on doing it.\footnote{Dave Smith, \emph{Libertarian National Convention 2024}, \url{https://rumble.com/v4yhk7e-may-30-2024.html}, t. 39:52}
\end{quote}

Yeah, I think my favourite part of Machiavelli's \emph{The Prince} is where he says that you have to trust what everybody tells you whilst you're competing with them for power, and then call them snakes and throw hissy fits should they betray you. These people, especially Rectenwald, spent months teasing the radical libertarians about how we don't understand politics and how they are the real gigachad understanders of the levers of power. They don't get to then turn around with this whole ``BUT ITS NOT FAIR! NOT FAIR NOT FAIR NOT FAIR! YOU PROMISED!!!'' whenever one of their political allies turns coat.

You also don't get to run this gambit of ``this guy is so dishonest, and he's a cop, and everybody should hate this guy for being an agent of the state'' bit when what you are pissed off at him for is that he didn't join \emph{your} club. If Ter Maat really is so terrible, why did you want him so bad?

Lastly, if Ter Maat really is to be painted as this evil guy because he went against his word--that is, lead the Mises Caucus to believe that he would endorse Rectenwald--why isn't Dave equally evil for leading the Mises Caucus to believe that he was planning to run, thus scarpering any hopes of them getting a good candidate, in their own words? Surely Dave's not running must have provided a much more substantial blow to the Mises Caucus' hopes than anything Ter Maat did; Dave seems to acknowledge this fact by stating that he takes the blame. But why is he not \emph{fully} taking the blame? That is, why is he running this bit of calling Ter Maat a liar and a cheat, when by his own standards he is a much bigger cheat? I know that Dave has only just gotten into the whole Christianity thing, but did he miss the verse about when one should cast the first stone?

Indeed, some good did come of the Mises Caucus' manoeuvring, namely that Ross was released. This is a massive victory for libertarianism, and easily the best thing the party has ever achieved. I agree that this alone is reason enough to justify going in with the Trump campaign. If I were able to vote in US elections, I probably would have voted for Trump on the basis of Ross Ulbricht and the potential of DOGE alone. But this is nowhere near the extent of everything that the Mises Caucus has done.

They have repeatedly, and without any self-reflection, bashed the radical basis of the movement, and attempted to sneak active socialists into the tent through the sally-port. Angela's error on this front is the active yellow-washing of MAGA Republicans she engaged in. Dave's error is actively endorsing the totalitarian policies that Trump advocated, differing only in that Dave thinks Trump hasn't gone far enough.
\section{The Philosophical Bust}
\label{sec:org9e07e94}
The whole border issue and the Mises Caucus' dreadful strategy are not entirely disconnected things as these people would have you believe,\footnote{Jacob Winograd tweet: "You're connecting two entirely different things/events
And also proving my point to @MonarchoHoppean that this is really just the fact that you guys are still peeved over a year later about Dave's comments
Let it go
Or don't 🤷🏻‍♂️" \url{https://x.com/BiblicalAnarchy/status/1908707891680264351} \url{https://archive.ph/DuXOK}} rather, they are both the direct result of the fundamental strategy that the neo-prags have taken up. All of it. The alienation of the radicals to court fiat libertarians to gain temporary culture war wins to eventually get---what?

Not libertarianism. I have already proven that Dave is not for liberty, let's take a review of the others. Angela sold out that label to Trump to claim a ``Libertarian POTUS,'' her husband constantly runs cover for all of Trumps disastrous and anti-libertarian policies, and Rectenwald is a professional leftist who was a true believer in Marxism right up until the machine he built turned on him. It is all the same strategy, the same strategy that required Dave to label us as ``worse than communists.'' We \emph{are} worse from his perspective, because we undermine him from within. We refute his entire plan. We negate the grift. Communists feed the grift.

\emph{That} is that they are really after. The grift. The grift is their categorial imperative and the driving force of everything they do. This is why they herd over libertarians in name only; great for grifting, terrible for winning. They are a poison that should be immediately excised. They are the anarcho-stirnerites and anarcho-pragmatists that Murray Rothbard warned us about. And per Rothbard's advice they should, on their own premises, shut up and get out of the fucking way of people trying to do the real work to bring about anarchy.

There is a phenomenon in economic theory called the business cycle. It is a procession of expansion in the economy, followed by a contraction; which goes on in a loop, every boom being followed by a bust. The Austrian account of this is that investors look to the amount saved to decide how long-run their investments should be. Consider the case of Robinson Crusoe on the desert island; if he has saved only a few twigs and logs, he should perhaps invest in a small raft, but if he has a vastly greater amount of capital saved he can build an entire fishing trawler, or cruise line.

The same is true of the regular economy. Investors gauge how many resources have been saved by looking at the interest rates of loans. If the interest rates are super low, this means that a whole lot of resources have been put aside, so investors can wait longer to be repaid. If the interest rates are high, this means that not many resources have been put aside, and accordingly investors will require a sooner payment on the loan.

When a fractional reserve bank prints off more money than actually exists, it signals to investors that more resources have been saved than have actually been saved. This ``fiat'' money has mislead investors into making longer-run investments than they otherwise would have, so when they find this out, it is suddenly revealed that all of these investments were wrong, and the market goes bust.

We can draw out of this economic case a more general principle. The economic boom was essentially caused by a friendly-sounding deception---a white lie if you will. Crusoe on our island is living in ignorance of the true state of his savings, because he has swallowed the falsehood that he is far better off than he really is. This falsehood will eventually run up against reality, and Crusoe will find that he has made incorrect decisions based upon his false premises. That's the general rule. When you operate upon false premises to generate some manner of ``fiat'' goodies--that is, goodies which don't actually exist--you will eventually reach some breaking point, where reality swoops in and demonstrates your choices to have been wrong. At such a point, a man has two options: he can either double-down and simply ignore this new fact of reality like he has been doing all along; or he can acknowledge the error of his ways and start to rebuild. The economic analogue for both would be either printing off more money to get another boom going, or letting the bubble burst, not bailing anyone out, and allowing capital to be reallocated to where it should have been in the first place.

Let us apply this generalised ``cycle of evasion'' to the case of the pragmatist strategy. Both the old prags and the neo-prags have as their guiding principle the herding over of fiat libertarians--that is, libertarians in name only--and the more they can herd into the tent the better they are said to be doing. This creates a libertarian ``boom'' period; where perhaps it becomes a very popular topic of discussion, and many people watch the podcasts, and many people attend the conferences. But, as discussed, this boom necessitates a bust. None of these people who were herded over were actually radicals for capitalism. None of them care about property rights as anything approaching a fundamental legal principle. They all inevitably wish to achieve some form of socialism in their various pet issues. It would be surprising if any one of them even understood the theory beyond the most basic soundbites.

After all, how could they ever advance beyond the stage of the soundbitatarian? These people are, by definition, followers; not leaders. This is why it is the same group of people who clap like seals as Dave pounds the drum in favour of principles, keep clapping like seals as he pounds the drum against principles. These people simply lack the conceptual framework to think for themselves.

By the fact that the pragmatist has taken as his guiding strategic light the herding in of sheep to generate fiat libertarians, he is selecting for people who do not carefully consider philosophy or come to their own conclusions on the various topics therein. Rather, he has selected for those who will just accept whatever hodgepodge of bundled philosophies which has been passed down to them; these bundles can come only in the form of soundbites.

By ``bundled philosophy'' I mean any philosophy which gives you a metaphysics (what is there?), epistemology (how do I know?), and ethics (what should I do about it?) in one neat bundle that can be passed from person to person. For instance, a particularly primitive form of religious proselytising might go like: what is there? God is there. How do I know? God told me. What should I do about it? Whatever God says.

Here the anarcho-pragmatist has focused his activities on providing a specific soundbite to slot into someone's ethics, without worrying about integrating it with any other area in philosophy. In fact, integration is to be discouraged, can't we all just get along and sing kumbaya!? This suits the pragmatist, his target demographic is not capable of integration in the first place. This is a far larger demographic to grift on, and it is far easier to just feed these people soundbites than to do the real work of philosophy.

Specifically, the anarcho-pragmatist will merely present some soundbite like ``socially liberal, fiscally conservative,'' or; ``don't hurt people, and don't take their stuff,'' or; ``black guns matter,'' or; ``did you guys know that war is bad?,'' or; ``we need to run government like a business!,'' and so on.

The pragmatist must ensure that the soundbites never get any more specific than this---being vague is the entire point. If some radical comes along and explains that ``don't hurt people, and don't take their stuff'' is to be laughed out of any serious discussion on jurisprudence and that the proper principle it is trying to point at is called the non-aggression principle, which is precisely defined in terms of praxeology; this man is to be denounced and laughed at as ``too autistic,'' or ``too idealistic,'' or ``living in anarcho-capitalism in his head.''

This is because if the radical is allowed to speak to the fiat libertarians, they will quickly be turned off, and be made to understand that they have no place in this movement. The radical will demonstrate to the fiatatarian that there are many parts of the tent he was herded into that he finds repulsive; many conclusions that are not so wholesome chungus and feminist or super duper trad and MAGA, depending, of course, on which variety of pragmatist did the herding. The fiatatarian will be disgusted by these radical conclusions and will quickly run away to be herded elsewhere at which point he may post on social media that he used to be a libertarian, before he grew up that is!

So what exactly is the practical implication of successful pragmatism---that is, pragmatism which generates many fiat libertarians. What is the consequence of the strategy which is revealed in the bust phase?

In the first case, the strategy involves a direct alienation of radicals. As discussed above, it is either-or. In the pursuit of milquetoast, pragmatist messaging, you have to simultaneously denounce radical messaging. One is forced to explain that they are not like those crazy anarchists, they just want reasonable limits on government overreach! The pragmatist must lionize the ``sensible'' or ``centrist'' path, and reject the ``idealistic'' or ``utopian'' one.

It is also not possible to have it both ways---to sit on the fence between radicalism and pragmatism. Those who attempt to ``be everyone's friend'' inevitably run up against this fact. In any opportunity where this type is given the choice between backing a radical or throwing him under the bus for a culture war win, they take the latter 10 times out of 10. If one of these kumbaya-atarians were to take the radical side, they become radicals and immediately place themselves in opposition to the pragmatists. In reality, kumbaya-atarianism is just another form of pragmatism; it is the denunciation of reality and principle, in favour of fiat and compromise. After all, can't we all just get along!? We cannot. There is no middle ground between A and non-A, the excluded middle does not exist. The kumbaya-atarian is a higher-order pragmatist; where the first order pragmatist attempts to bridge between radicals and normies by herding the latter into the libertarian tent, the second order pragmatist attempts to bridge between pragmatists and radicals by herding the pragmatists into the radical tent. An attempt which fails for obvious reasons.

Not only does the pragmatist strategy necessitate the direct alienation of radicals, it also indirectly alienates them by making coordination between radicals more difficult. If I find myself in this tent which is being constantly fluffed up by an influx of people who hate libertarianism, but have had the label hastily stamped on their forehead upon entry, it becomes increasingly more challenging to find other radicals in the movement to converse and organise with.

One might find himself in a supposedly-libertarian party or conference and be faced with legion subversions to the theory. ``Oh, I am definitely a libertarian, but it violates my NAP to allow any immigration;'' ``oh, you misunderstand, it is actually an initiation of force to lie about someone as it damages their reputation: free speech in moderation, I say!;'' ``oh, the non-aggression principle is definitely true, but aggression means when the Proletariat are exploited, so we need socialism to be truly free.''

These subversions have two sources: (1) through fiatatarians bringing along their old bad premises, but translating them into libertarian terminology; and (2) through active subverters in the movement. These subverters come into the tent through the gaping hole made by the pragmatists, seek out new and easily swayed advocates of liberty, and fill their heads with poison. These subverters are in the business of trying to herd people out of the liberty tent, and into their own ideology's tent. Even if they fail to lead the individual completely back home, they often succeed in demonstrating to him that the libertarian movement is full of bad actors, and is not to be engaged with. This counts as victory also.

All of this indirect alienation ends up in the same place: the radical vanguard is slowed down, and slowly loses members to attrition. Potential members of the vanguard are unable to find it, and are left wandering forever. And without this base of radicals the movement becomes doomed to failure, perhaps needing to be spun off into a new ideology, leaving the old one to rot.

The movement \emph{must} die without the radicals because without the radicals there is not, strictly speaking, any movement at all. That is, the radical base are not merely the core, or guiding light, or misguided autistic wing of the movement; the radicals \emph{are} the movement. Every single other person in the tent is either entirely irrelevant or an active hindrance to the movement. Properly speaking, one cannot describe this ballast as part of any movement as they generate no \emph{motion}. They exist only as inertia, they go along in the same direction unless acted upon by some force generated by radicals. Having this inertia is a very bad thing for a small ideology, because the active subverters can provide this force just as fully as the radical vanguard.

Very simply, my prescription is thus: under no circumstances must you attempt to do any herding of sheep people into the libertarian tent. Simply state the case in its most radical form, and smile as the majority peel away to do nothing of substance. That is how we win.
\section{Smith Contra Smith}
\label{sec:orged01f03}
There are a great many disagreements among those who call themselves libertarians. As the old joke goes, if you have three libertarians in a room, the only thing they will agree on is that there is only one libertarian in that room. Of course, there is an objective standard of what counts as libertarianism, consisting of a strict philosophical adherence to the non-aggression principle, but I want to step outside of that for just a moment. I want to demonstrate that the great villain of this video, Dave Smith, is not a libertarian--not by my standard--but by his own.

Dave tells us what a libertarian is during an interview with Reason magazine:
\begin{quote}
Liz Wolfe: What is a libertarian? What beliefs are disqualifying for libertarians to hold today?

Dave Smith: [\ldots{}] libertarianism is the belief in self-ownership, private property rights, and the non-aggression principle. That is the best, philosophically-sound definition of it that isn't circular.\footnote{ReasonTV, \emph{What is a libertarian? | Dave Smith | Just Asking Questions - Ep. 1}, \url{https://www.youtube.com/watch?v=SByjfbOUvWc}, t. 1:10}
\end{quote}

Ok, so Dave's criteria for being a libertarian are that you must believe in:
\begin{enumerate}
\item self-ownership;
\item property rights, and;
\item the non-aggression principle.
\end{enumerate}

Now, we can actually narrow this down into a single criterion. You see, the non-aggression principle implies \emph{both} private property rights and self-ownership. It implies self-ownership, because anybody who comes and tries to use your body for something you don't want to use it for is necessarily initiating a conflict with you, because you have the best objective link to your body. This means that you own yourself. You can have private property rights in external means for the same reason; if you are the first-comer to some good, it is necessarily the case that some later-comer is initiating a conflict with you if they try to take it off of you. So by Dave's standard, what it means to be a libertarian is to have a philosophical adherence to the non-aggression principle.

That is, a libertarian is someone who in any matter that law tackles, will adhere to the non-aggression principle above all else. A libertarian might think that it is a really bad idea to carry firearms, or think that adults should not voluntarily agree to sex re-assignment surgery; but the libertarian would discount these irrelevant ethical principles as far as the question of rights is concerned. So let's see whether Dave is able to meet this criterion.

\begin{quote}
You know there's--I'm not strawmanning--but there are some libertarians, not serious libertarians, or prominent libertarians, but there's enough of the rank-and-file member who will actually say, ``well, we should go further, and in fact all of these people should be allowed in government schools, or government libraries, or government playgrounds, or things like that.''

To me, that is incorrect by libertarian theory--I want to make that clear first--but I will also say: if that was correct by libertarian theory, then I would simply reject libertarianism. If the case was that random homeless drug addicts who are covered in their own faeces and filth should be allowed in public schools, [\ldots{}] then I would just abandon libertarianism.\footnote{The Biblical Anarchy Podcast, \emph{What is the Libertarian Response to Immigration w/Dave Smith}, \url{https://www.youtube.com/watch?v=ghqN\_F2iTdU}, 13:15}
\end{quote}

Notice what the question here is: the question is whether random people--whether they be homeless drug addicts or whatever--have the right to enter various possessions that are being forestalled by the government. The examples given being public schools, libraries, and playgrounds. I have already gone into this issue in depth earlier in the video, so I will not re-iterate that here; rather I want to point to how \emph{Dave} is responding to this. He asserts that on libertarian theory, those people have no such right, but he doesn't stop there. He then goes on to say that if someone were to present him a proof that given the NAP, those people in fact do have that right, then he would abandon libertarianism. That is, he would drop the NAP from his philosophy before he would drop the ``no homeless in public parks'' principle, or whatever the hell his general principle is on that front.

In no sense can Dave be said to have a philosophical adherence to the NAP in this case. Again, what a libertarian would mean--\emph{by Dave's own standard}--is to use the NAP as the starting point in any analysis on questions that law tackles. If Dave was using the NAP as a starting point, he could never come to a question in law where he throws away the NAP for some ill-defined, vibes-based principle, like he does on the homeless question.

It is specifically the statists who come to some question about the rights of the matter; find that a right held by X person would yield behaviour they find untoward, and then conclude that therefore no such right exists. That is an inversion of the libertarian way of dealing with the law. The libertarian starts from the premise of the NAP, from there he \emph{discovers} what rights people have, and that is the end of the discussion as far as law is concerned.

Dave has just admitted to everyone listening that he has adopted a fundamentally unlibertarian theory of rights. Dave has certain conclusions about what society should look like, and he has only happened to work back from those conclusions to come to the NAP \emph{so far}. If at any point an Andrew Wilson were to come along and prove to Dave that his preferred bureaucratic structure actually requires some alternative theory of rights, Dave would drop the NAP.

Dave Smith is not a libertarian, so says Dave Smith.

The pragmatist subversion is not the only counter to anarcho-capitalism present today; there are innumerable others. If you wish to know the hardcore ancap case and how it counters these other subversions you have to watch this video where I respond to counterarguments from TIKHistory, Ayn Rand, and everyone else I could find.
\end{document}
